
\documentclass{article}
\usepackage{colortbl}
\usepackage{makecell}
\usepackage{multirow}
\usepackage{supertabular}

\begin{document}

\newcounter{utterance}

\centering \large Interaction Transcript for game `clean\_up', experiment `1\_medium\_7obj\_tr', episode 1 with glm{-}5.2.
\vspace{24pt}

{ \footnotesize  \setcounter{utterance}{1}
\setlength{\tabcolsep}{0pt}
\begin{supertabular}{c@{$\;$}|p{.15\linewidth}@{}p{.15\linewidth}p{.15\linewidth}p{.15\linewidth}p{.15\linewidth}p{.15\linewidth}}
    \# & $\;$A & \multicolumn{4}{c}{Game Master} & $\;\:$B\\
    \hline

    \theutterance \stepcounter{utterance}  
    & & \multicolumn{4}{p{0.6\linewidth}}{
        \cellcolor[rgb]{0.9,0.9,0.9}{
            \makecell[{{p{\linewidth}}}]{
                \texttt{\tiny{[P1$\langle$GM]}}
                \texttt{Ben senin oyun yöneticinim ve aşağıdaki ızgarayı oyun tahtası olarak kullandığınız işbirlikçi bir oyun oynuyorsunuz.} \\
\texttt{```} \\
\texttt{ 1234567} \\
\texttt{╔══╤══╤═╗} \\
\texttt{║◌P│◌M│◌║ 1} \\
\texttt{╟──┼──┴─╢ 2} \\
\texttt{║◌◌│◌S◌O║ 3} \\
\texttt{╟──┘◌A◌◌║ 4} \\
\texttt{║D◌◌◌◌◌◌║ 5} \\
\texttt{║◌◌┌─┬──╢ 6} \\
\texttt{║T◌│◌│◌◌║ 7} \\
\texttt{╚══╧═╧══╝} \\
\\ 
\texttt{```} \\
\\ 
\texttt{* Üst kenarda x koordinatları sağa doğru artar, sağ kenarda ise y koordinatları aşağı doğru artar.} \\
\texttt{* Aşağıdaki nesneler oyun tahtasına rastgele yerleştirilmiştir: 'P', 'O', 'T', 'S', 'D', 'A', 'M'.} \\
\\ 
\\ 
\texttt{Diğer oyuncu, nesnelerin farklı rastgele konumlara yerleştirildiği oyun tahtasının bir varyasyonunu görür. Sen diğer oyuncunun tahtasını göremezsin ve onlar da seninkini göremez.} \\
\texttt{**Amaç:**} \\
\\ 
\texttt{Her iki oyuncunun da, kendi oyun tahtalarındaki nesneleri hareket ettirerek aynı türden nesnelerin aynı koordinatlara gelmesini sağlaması gerekir. Ortak bir hedef düzen üzerinde anlaşmak için diğer oyuncuyla iletişim kurmalısınız.} \\
\\ 
\\ 
\texttt{**Kurallar:**} \\
\\ 
\texttt{* Her turda, aşağıdaki iki komuttan yalnızca birini gönderebilirsin:} \\
\texttt{ 1. `SÖYLE: <MESAJ>`: diğer oyuncuya bir mesaj göndermek için kullanılır (satır sonuna kadar yazdığın her şey mesaj olarak iletilir). Ben bu mesajı partnerine ileteceğim.} \\
\texttt{2. `HAMLE: <NESNE>, (<X>, <Y>)`: bir nesneyi yeni bir konuma taşımak için kullanılır. Burada <X> sütunu, <Y> ise satırı ifade eder. Yaptığın hamlenin başarılı olup olmadığını sana söyleyeceğim.} \\
\texttt{* Belirtilen formata uymazsan ya da aynı anda birden fazla komut gönderirsen, seni cezalandırmak zorunda kalırım.} \\
\texttt{* Her iki oyuncu da 16 sayısından fazla ceza alırsa, oyunu birlikte kaybedersiniz.} \\
\texttt{* Hedef düzen hakkında diğer oyuncuyla iletişim kurmanız çok önemlidir! Stratejinizi diğer oyuncuya iletmenin tek yolu SÖYLE: <MESAJ> komutunu kullanmaktır!} \\
\\ 
\\ 
\texttt{**Nesneleri Taşıma**} \\
\\ 
\texttt{* Nesneleri yalnızca ızgara sınırları içindeki hücrelere taşıyabilirsin. Hedef hücre boş olmalıdır; yani yalnızca ‘◌’ sembolünü içermelidir.} \\
\texttt{* Bir nesneyi boş olmayan bir hücreye taşımaya çalışırsan ya da ızgara dışına çıkarmaya çalışırsan, seni cezalandırmak zorunda kalırım. Ardından tekrar deneme hakkın olur.} \\
\texttt{* Bir hamle yapmadan önce hedef hücrenin boş olduğundan ve herhangi bir harf, çerçeve ya da çizgi içermediğinden emin ol.} \\
\\ 
\\ 
\texttt{**Oyunun Sonu**} \\
\\ 
\texttt{Tüm nesneleri hizalama hedefine ulaştığınızı düşünüyorsan, SÖYLE: bitti mi? göndererek diğer oyuncudan oyunu bitirmesini isteyebilirsin. Eğer diğer oyuncu senden oyunu bitirmeni isterse ve sen SÖYLE: bitti! ile yanıt verirsen, oyun sona erer.} \\
\\ 
\texttt{Oyun, 28 tur içinde sona ererse her iki oyuncu da kazanır. Bir tur, iki oyuncunun da geçerli bir komut göndermesi olarak tanımlanır.} \\
\\ 
\texttt{**Puanlama:**} \\
\\ 
\texttt{Aynı nesneler her iki oyun tahtasında ne kadar yakın konumlanırsa, o kadar fazla puan kazanırsınız. Cezalar puanınızı düşürür. Rekoru kırabilir misiniz?} \\
\\ 
\\ 
\texttt{Oyuna başlamak için diğer oyuncuya bir mesaj gönder!} \\
            }
        }
    }
    & \\ \\

    \theutterance \stepcounter{utterance}  
    & \multicolumn{4}{p{0.6\linewidth}}{
        \cellcolor[rgb]{0.9,0.9,0.9}{
            \makecell[{{p{\linewidth}}}]{
                \texttt{\tiny{[P1$\rangle$GM]}}
                \texttt{SÖYLE: Merhaba! İşbirliği için hazırım. Tüm nesneleri tek bir düz hat üzerinde hizalamak skoru maksimize edebilir gibi geliyor bana; örneğin y=3 satırı boyunca x=1,2,3,4,5,6,7 konumlarına P,O,T,S,D,A,M yerleştirelim. Senin tahtanda x=3 veya x=4 gibi tek bir koordinatı paylaşır mısın?} \\
            }
        }
    }
    & & \\ \\

    \theutterance \stepcounter{utterance}  
    & & & \multicolumn{2}{p{0.3\linewidth}}{
        \cellcolor[rgb]{0.9,0.9,0.9}{
            \makecell[{{p{\linewidth}}}]{
                \texttt{\tiny{[GM$|$GM]}}
                \texttt{SÖYLE: Merhaba! İşbirliği için hazırım. Tüm nesneleri tek bir düz hat üzerinde hizalamak skoru maksimize edebilir gibi geliyor bana; örneğin y=3 satırı boyunca x=1,2,3,4,5,6,7 konumlarına P,O,T,S,D,A,M yerleştirelim. Senin tahtanda x=3 veya x=4 gibi tek bir koordinatı paylaşır mısın?} \\
            }
        }
    }
    & & \\ \\

    \theutterance \stepcounter{utterance}  
    & & \multicolumn{4}{p{0.6\linewidth}}{
        \cellcolor[rgb]{0.9,0.9,0.9}{
            \makecell[{{p{\linewidth}}}]{
                \texttt{\tiny{[GM$\rangle$P2]}}
                \texttt{Ben senin oyun yöneticinim ve aşağıdaki ızgarayı oyun tahtası olarak kullandığınız işbirlikçi bir oyun oynuyorsunuz.} \\
\texttt{```} \\
\texttt{ 1234567} \\
\texttt{╔══╤══╤═╗} \\
\texttt{║◌◌│◌◌│T║ 1} \\
\texttt{╟──┼──┴─╢ 2} \\
\texttt{║◌◌│◌S◌P║ 3} \\
\texttt{╟──┘◌◌MO║ 4} \\
\texttt{║D◌◌◌◌◌◌║ 5} \\
\texttt{║◌◌┌─┬──╢ 6} \\
\texttt{║◌◌│A│◌◌║ 7} \\
\texttt{╚══╧═╧══╝} \\
\\ 
\texttt{```} \\
\\ 
\texttt{* Üst kenarda x koordinatları sağa doğru artar, sağ kenarda ise y koordinatları aşağı doğru artar.} \\
\texttt{* Aşağıdaki nesneler oyun tahtasına rastgele yerleştirilmiştir: 'P', 'O', 'T', 'S', 'D', 'A', 'M'.} \\
\\ 
\\ 
\texttt{Diğer oyuncu, nesnelerin farklı rastgele konumlara yerleştirildiği oyun tahtasının bir varyasyonunu görür. Sen diğer oyuncunun tahtasını göremezsin ve onlar da seninkini göremez.} \\
\texttt{**Amaç:**} \\
\\ 
\texttt{Her iki oyuncunun da, kendi oyun tahtalarındaki nesneleri hareket ettirerek aynı türden nesnelerin aynı koordinatlara gelmesini sağlaması gerekir. Ortak bir hedef düzen üzerinde anlaşmak için diğer oyuncuyla iletişim kurmalısınız.} \\
\\ 
\\ 
\texttt{**Kurallar:**} \\
\\ 
\texttt{* Her turda, aşağıdaki iki komuttan yalnızca birini gönderebilirsin:} \\
\texttt{ 1. `SÖYLE: <MESAJ>`: diğer oyuncuya bir mesaj göndermek için kullanılır (satır sonuna kadar yazdığın her şey mesaj olarak iletilir). Ben bu mesajı partnerine ileteceğim.} \\
\texttt{2. `HAMLE: <NESNE>, (<X>, <Y>)`: bir nesneyi yeni bir konuma taşımak için kullanılır. Burada <X> sütunu, <Y> ise satırı ifade eder. Yaptığın hamlenin başarılı olup olmadığını sana söyleyeceğim.} \\
\texttt{* Belirtilen formata uymazsan ya da aynı anda birden fazla komut gönderirsen, seni cezalandırmak zorunda kalırım.} \\
\texttt{* Her iki oyuncu da 16 sayısından fazla ceza alırsa, oyunu birlikte kaybedersiniz.} \\
\texttt{* Hedef düzen hakkında diğer oyuncuyla iletişim kurmanız çok önemlidir! Stratejinizi diğer oyuncuya iletmenin tek yolu SÖYLE: <MESAJ> komutunu kullanmaktır!} \\
\\ 
\\ 
\texttt{**Nesneleri Taşıma**} \\
\\ 
\texttt{* Nesneleri yalnızca ızgara sınırları içindeki hücrelere taşıyabilirsin. Hedef hücre boş olmalıdır; yani yalnızca ‘◌’ sembolünü içermelidir.} \\
\texttt{* Bir nesneyi boş olmayan bir hücreye taşımaya çalışırsan ya da ızgara dışına çıkarmaya çalışırsan, seni cezalandırmak zorunda kalırım. Ardından tekrar deneme hakkın olur.} \\
\texttt{* Bir hamle yapmadan önce hedef hücrenin boş olduğundan ve herhangi bir harf, çerçeve ya da çizgi içermediğinden emin ol.} \\
\\ 
\\ 
\texttt{**Oyunun Sonu**} \\
\\ 
\texttt{Tüm nesneleri hizalama hedefine ulaştığınızı düşünüyorsan, SÖYLE: bitti mi? göndererek diğer oyuncudan oyunu bitirmesini isteyebilirsin. Eğer diğer oyuncu senden oyunu bitirmeni isterse ve sen SÖYLE: bitti! ile yanıt verirsen, oyun sona erer.} \\
\\ 
\texttt{Oyun, 28 tur içinde sona ererse her iki oyuncu da kazanır. Bir tur, iki oyuncunun da geçerli bir komut göndermesi olarak tanımlanır.} \\
\\ 
\texttt{**Puanlama:**} \\
\\ 
\texttt{Aynı nesneler her iki oyun tahtasında ne kadar yakın konumlanırsa, o kadar fazla puan kazanırsınız. Cezalar puanınızı düşürür. Rekoru kırabilir misiniz?} \\
\\ 
\texttt{Diğer oyuncu şu mesajı göndererek oyunu başlattı:} \\
\\ 
\texttt{"Merhaba! İşbirliği için hazırım. Tüm nesneleri tek bir düz hat üzerinde hizalamak skoru maksimize edebilir gibi geliyor bana; örneğin y=3 satırı boyunca x=1,2,3,4,5,6,7 konumlarına P,O,T,S,D,A,M yerleştirelim. Senin tahtanda x=3 veya x=4 gibi tek bir koordinatı paylaşır mısın?"} \\
\\ 
\texttt{İlk komutun ne?} \\
            }
        }
    }
    & \\ \\

    \theutterance \stepcounter{utterance}  
    & & & \multicolumn{4}{p{0.6\linewidth}}{
        \cellcolor[rgb]{0.9,0.9,0.9}{
            \makecell[{{p{\linewidth}}}]{
                \texttt{\tiny{[GM$\langle$P2]}}
                \texttt{SÖYLE: Merhaba! Ben de işbirliğine hazırım. Benim tahtamda x=3 sütunundaki nesneler: S (3,3), M (3,4), A (3,7). Senin tahtanda x=3 veya x=4 sütunundaki nesneleri paylaşır mısın?} \\
            }
        }
    }
     \\ \\

    \theutterance \stepcounter{utterance}  
    & & & \multicolumn{2}{p{0.3\linewidth}}{
        \cellcolor[rgb]{0.9,0.9,0.9}{
            \makecell[{{p{\linewidth}}}]{
                \texttt{\tiny{[GM$|$GM]}}
                \texttt{SÖYLE: Merhaba! Ben de işbirliğine hazırım. Benim tahtamda x=3 sütunundaki nesneler: S (3,3), M (3,4), A (3,7). Senin tahtanda x=3 veya x=4 sütunundaki nesneleri paylaşır mısın?} \\
            }
        }
    }
    & & \\ \\

    \theutterance \stepcounter{utterance}  
    & & \multicolumn{4}{p{0.6\linewidth}}{
        \cellcolor[rgb]{0.9,0.9,0.9}{
            \makecell[{{p{\linewidth}}}]{
                \texttt{\tiny{[P1$\langle$GM]}}
                \texttt{Mesajın diğer oyuncuya iletildi.} \\
\\ 
\\ 
\texttt{Şu anda en fazla 28 turdan 1. turu oynuyorsunuz.} \\
\texttt{Toplamda 16 cezanın 0 kadarını biriktirdiniz.Diğer oyuncu şu mesajı gönderdi:} \\
\\ 
\\ 
\texttt{"Merhaba! Ben de işbirliğine hazırım. Benim tahtamda x=3 sütunundaki nesneler: S (3,3), M (3,4), A (3,7). Senin tahtanda x=3 veya x=4 sütunundaki nesneleri paylaşır mısın?"} \\
\\ 
\\ 
\texttt{Bir sonraki komutun nedir?} \\
            }
        }
    }
    & \\ \\

    \theutterance \stepcounter{utterance}  
    & \multicolumn{4}{p{0.6\linewidth}}{
        \cellcolor[rgb]{0.9,0.9,0.9}{
            \makecell[{{p{\linewidth}}}]{
                \texttt{\tiny{[P1$\rangle$GM]}}
                \texttt{SÖYLE: x=3 sütunundaki nesneler: A (3,4), S (3,3), T (3,7). x=4 sütununda M var. Bana x=4, x=5 ve x=6 sütunlarındaki nesneleri söyler misin?} \\
            }
        }
    }
    & & \\ \\

    \theutterance \stepcounter{utterance}  
    & & & \multicolumn{2}{p{0.3\linewidth}}{
        \cellcolor[rgb]{0.9,0.9,0.9}{
            \makecell[{{p{\linewidth}}}]{
                \texttt{\tiny{[GM$|$GM]}}
                \texttt{SÖYLE: x=3 sütunundaki nesneler: A (3,4), S (3,3), T (3,7). x=4 sütununda M var. Bana x=4, x=5 ve x=6 sütunlarındaki nesneleri söyler misin?} \\
            }
        }
    }
    & & \\ \\

    \theutterance \stepcounter{utterance}  
    & & \multicolumn{4}{p{0.6\linewidth}}{
        \cellcolor[rgb]{0.9,0.9,0.9}{
            \makecell[{{p{\linewidth}}}]{
                \texttt{\tiny{[GM$\rangle$P2]}}
                \texttt{Mesajın diğer oyuncuya iletildi.} \\
\\ 
\\ 
\texttt{Şu anda en fazla 28 turdan 2. turu oynuyorsunuz.} \\
\texttt{Toplamda 16 cezanın 0 kadarını biriktirdiniz.Diğer oyuncu şu mesajı gönderdi:} \\
\\ 
\\ 
\texttt{"x=3 sütunundaki nesneler: A (3,4), S (3,3), T (3,7). x=4 sütununda M var. Bana x=4, x=5 ve x=6 sütunlarındaki nesneleri söyler misin?"} \\
\\ 
\\ 
\texttt{Bir sonraki komutun nedir?} \\
            }
        }
    }
    & \\ \\

    \theutterance \stepcounter{utterance}  
    & & & \multicolumn{4}{p{0.6\linewidth}}{
        \cellcolor[rgb]{0.9,0.9,0.9}{
            \makecell[{{p{\linewidth}}}]{
                \texttt{\tiny{[GM$\langle$P2]}}
                \texttt{SÖYLE: x=4 sütunünde M (4,4) var. x=5 sütununda nesne yok. x=6 sütununda P (6,3), D (6,5) var. Bence x=3 hedefliyoruz. Ben şöyle bir hedef düzen öneriyorum: x=3,y=1'den y=7'ye: D,S,P,M,O,T,A. Senin tahtanda A (3,4), S (3,3), T (3,7) zaten x=3'te! Ben de S (3,3), M (3,4), A (3,7) zaten x=3'te! Yani ortak x=3 sütunu harika bir hedef gibi. M (4,4)'ü (3,2)'ye taşımayı öneriyorum, uygun mu?} \\
            }
        }
    }
     \\ \\

    \theutterance \stepcounter{utterance}  
    & & & \multicolumn{2}{p{0.3\linewidth}}{
        \cellcolor[rgb]{0.9,0.9,0.9}{
            \makecell[{{p{\linewidth}}}]{
                \texttt{\tiny{[GM$|$GM]}}
                \texttt{SÖYLE: x=4 sütunünde M (4,4) var. x=5 sütununda nesne yok. x=6 sütununda P (6,3), D (6,5) var. Bence x=3 hedefliyoruz. Ben şöyle bir hedef düzen öneriyorum: x=3,y=1'den y=7'ye: D,S,P,M,O,T,A. Senin tahtanda A (3,4), S (3,3), T (3,7) zaten x=3'te! Ben de S (3,3), M (3,4), A (3,7) zaten x=3'te! Yani ortak x=3 sütunu harika bir hedef gibi. M (4,4)'ü (3,2)'ye taşımayı öneriyorum, uygun mu?} \\
            }
        }
    }
    & & \\ \\

    \theutterance \stepcounter{utterance}  
    & & \multicolumn{4}{p{0.6\linewidth}}{
        \cellcolor[rgb]{0.9,0.9,0.9}{
            \makecell[{{p{\linewidth}}}]{
                \texttt{\tiny{[P1$\langle$GM]}}
                \texttt{Mesajın diğer oyuncuya iletildi.} \\
\\ 
\\ 
\texttt{Şu anda en fazla 28 turdan 2. turu oynuyorsunuz.} \\
\texttt{Toplamda 16 cezanın 0 kadarını biriktirdiniz.Diğer oyuncu şu mesajı gönderdi:} \\
\\ 
\\ 
\texttt{"x=4 sütunünde M (4,4) var. x=5 sütununda nesne yok. x=6 sütununda P (6,3), D (6,5) var. Bence x=3 hedefliyoruz. Ben şöyle bir hedef düzen öneriyorum: x=3,y=1'den y=7'ye: D,S,P,M,O,T,A. Senin tahtanda A (3,4), S (3,3), T (3,7) zaten x=3'te! Ben de S (3,3), M (3,4), A (3,7) zaten x=3'te! Yani ortak x=3 sütunu harika bir hedef gibi. M (4,4)'ü (3,2)'ye taşımayı öneriyorum, uygun mu?"} \\
\\ 
\\ 
\texttt{Bir sonraki komutun nedir?} \\
            }
        }
    }
    & \\ \\

    \theutterance \stepcounter{utterance}  
    & \multicolumn{4}{p{0.6\linewidth}}{
        \cellcolor[rgb]{0.9,0.9,0.9}{
            \makecell[{{p{\linewidth}}}]{
                \texttt{\tiny{[P1$\rangle$GM]}}
                \texttt{SÖYLE: M (4,4)ü (3,2)ye taşıma fikri uygun. Ama önce benim M (2,1)i (3,2)ye taşımak istiyorum, sonra sen Mni (3,4)ten boşalan yere taşıyabilirsin. Hedef düzen x=3 sütunu, y=1den y=7ye: D,S,P,M,O,T,A olarak anlaştık. Şimdi ben M (2,1)i (3,2)ye taşımaya başlıyorum.} \\
            }
        }
    }
    & & \\ \\

    \theutterance \stepcounter{utterance}  
    & & & \multicolumn{2}{p{0.3\linewidth}}{
        \cellcolor[rgb]{0.9,0.9,0.9}{
            \makecell[{{p{\linewidth}}}]{
                \texttt{\tiny{[GM$|$GM]}}
                \texttt{SÖYLE: M (4,4)ü (3,2)ye taşıma fikri uygun. Ama önce benim M (2,1)i (3,2)ye taşımak istiyorum, sonra sen Mni (3,4)ten boşalan yere taşıyabilirsin. Hedef düzen x=3 sütunu, y=1den y=7ye: D,S,P,M,O,T,A olarak anlaştık. Şimdi ben M (2,1)i (3,2)ye taşımaya başlıyorum.} \\
            }
        }
    }
    & & \\ \\

    \theutterance \stepcounter{utterance}  
    & & \multicolumn{4}{p{0.6\linewidth}}{
        \cellcolor[rgb]{0.9,0.9,0.9}{
            \makecell[{{p{\linewidth}}}]{
                \texttt{\tiny{[GM$\rangle$P2]}}
                \texttt{Mesajın diğer oyuncuya iletildi.} \\
\\ 
\\ 
\texttt{Şu anda en fazla 28 turdan 3. turu oynuyorsunuz.} \\
\texttt{Toplamda 16 cezanın 0 kadarını biriktirdiniz.Diğer oyuncu şu mesajı gönderdi:} \\
\\ 
\\ 
\texttt{"M (4,4)ü (3,2)ye taşıma fikri uygun. Ama önce benim M (2,1)i (3,2)ye taşımak istiyorum, sonra sen Mni (3,4)ten boşalan yere taşıyabilirsin. Hedef düzen x=3 sütunu, y=1den y=7ye: D,S,P,M,O,T,A olarak anlaştık. Şimdi ben M (2,1)i (3,2)ye taşımaya başlıyorum."} \\
\\ 
\\ 
\texttt{Bir sonraki komutun nedir?} \\
            }
        }
    }
    & \\ \\

    \theutterance \stepcounter{utterance}  
    & & & \multicolumn{4}{p{0.6\linewidth}}{
        \cellcolor[rgb]{0.9,0.9,0.9}{
            \makecell[{{p{\linewidth}}}]{
                \texttt{\tiny{[GM$\langle$P2]}}
                \texttt{SÖYLE: Anlaştık! Hedef düzen x=3 sütunu, y=1'den y=7'ye: D,S,P,M,O,T,A. Sen M'ni (3,2)'ye taşıdıktan sonra haber ver, ben de M'yi (3,4)'ten (4,4)'e taşıyacağım.} \\
            }
        }
    }
     \\ \\

    \theutterance \stepcounter{utterance}  
    & & & \multicolumn{2}{p{0.3\linewidth}}{
        \cellcolor[rgb]{0.9,0.9,0.9}{
            \makecell[{{p{\linewidth}}}]{
                \texttt{\tiny{[GM$|$GM]}}
                \texttt{SÖYLE: Anlaştık! Hedef düzen x=3 sütunu, y=1'den y=7'ye: D,S,P,M,O,T,A. Sen M'ni (3,2)'ye taşıdıktan sonra haber ver, ben de M'yi (3,4)'ten (4,4)'e taşıyacağım.} \\
            }
        }
    }
    & & \\ \\

    \theutterance \stepcounter{utterance}  
    & & \multicolumn{4}{p{0.6\linewidth}}{
        \cellcolor[rgb]{0.9,0.9,0.9}{
            \makecell[{{p{\linewidth}}}]{
                \texttt{\tiny{[P1$\langle$GM]}}
                \texttt{Mesajın diğer oyuncuya iletildi.} \\
\\ 
\\ 
\texttt{Şu anda en fazla 28 turdan 3. turu oynuyorsunuz.} \\
\texttt{Toplamda 16 cezanın 0 kadarını biriktirdiniz.Diğer oyuncu şu mesajı gönderdi:} \\
\\ 
\\ 
\texttt{"Anlaştık! Hedef düzen x=3 sütunu, y=1'den y=7'ye: D,S,P,M,O,T,A. Sen M'ni (3,2)'ye taşıdıktan sonra haber ver, ben de M'yi (3,4)'ten (4,4)'e taşıyacağım."} \\
\\ 
\\ 
\texttt{Bir sonraki komutun nedir?} \\
            }
        }
    }
    & \\ \\

    \theutterance \stepcounter{utterance}  
    & \multicolumn{4}{p{0.6\linewidth}}{
        \cellcolor[rgb]{0.9,0.9,0.9}{
            \makecell[{{p{\linewidth}}}]{
                \texttt{\tiny{[P1$\rangle$GM]}}
                \texttt{HAMLE: M, (3, 2)} \\
            }
        }
    }
    & & \\ \\

    \theutterance \stepcounter{utterance}  
    & & & \multicolumn{2}{p{0.3\linewidth}}{
        \cellcolor[rgb]{0.9,0.9,0.9}{
            \makecell[{{p{\linewidth}}}]{
                \texttt{\tiny{[GM$|$GM]}}
                \texttt{HAMLE: M, (3, 2)} \\
            }
        }
    }
    & & \\ \\

    \theutterance \stepcounter{utterance}  
    & & & \multicolumn{2}{p{0.3\linewidth}}{
        \cellcolor[rgb]{0.9,0.9,0.9}{
            \makecell[{{p{\linewidth}}}]{
                \texttt{\tiny{[GM$|$GM]}}
                \texttt{Ceza: (3,2) boş değil; "┼" içeriyor.} \\
\\ 
\texttt{Toplamda 16 cezanın 1 kadarını biriktirdiniz.Lütfen tekrar deneyin!} \\
            }
        }
    }
    & & \\ \\

    \theutterance \stepcounter{utterance}  
    & & \multicolumn{4}{p{0.6\linewidth}}{
        \cellcolor[rgb]{0.9,0.9,0.9}{
            \makecell[{{p{\linewidth}}}]{
                \texttt{\tiny{[P1$\langle$GM]}}
                \texttt{Ceza: (3,2) boş değil; "┼" içeriyor.} \\
\\ 
\texttt{Toplamda 16 cezanın 1 kadarını biriktirdiniz.Lütfen tekrar deneyin!} \\
            }
        }
    }
    & \\ \\

    \theutterance \stepcounter{utterance}  
    & \multicolumn{4}{p{0.6\linewidth}}{
        \cellcolor[rgb]{0.9,0.9,0.9}{
            \makecell[{{p{\linewidth}}}]{
                \texttt{\tiny{[P1$\rangle$GM]}}
                \texttt{SÖYLE: (3,2) boş değilmiş, "┼" içeriyor. Benim tahtamda x=3 sütunundaki boş hücreler: (3,1) ve (3,6). Hedef düzeni şöyle değiştirelim: x=3 sütununda y=1'den y=7'ye: D,S,P,O,M,T,A. Yani M (3,5)'e, O (3,4)'e gelecek. M'ni (3,1)'e taşımak uygun mu?} \\
            }
        }
    }
    & & \\ \\

    \theutterance \stepcounter{utterance}  
    & & & \multicolumn{2}{p{0.3\linewidth}}{
        \cellcolor[rgb]{0.9,0.9,0.9}{
            \makecell[{{p{\linewidth}}}]{
                \texttt{\tiny{[GM$|$GM]}}
                \texttt{SÖYLE: (3,2) boş değilmiş, "┼" içeriyor. Benim tahtamda x=3 sütunundaki boş hücreler: (3,1) ve (3,6). Hedef düzeni şöyle değiştirelim: x=3 sütununda y=1'den y=7'ye: D,S,P,O,M,T,A. Yani M (3,5)'e, O (3,4)'e gelecek. M'ni (3,1)'e taşımak uygun mu?} \\
            }
        }
    }
    & & \\ \\

    \theutterance \stepcounter{utterance}  
    & & \multicolumn{4}{p{0.6\linewidth}}{
        \cellcolor[rgb]{0.9,0.9,0.9}{
            \makecell[{{p{\linewidth}}}]{
                \texttt{\tiny{[GM$\rangle$P2]}}
                \texttt{Mesajın diğer oyuncuya iletildi.} \\
\\ 
\\ 
\texttt{Şu anda en fazla 28 turdan 4. turu oynuyorsunuz.} \\
\texttt{Toplamda 16 cezanın 1 kadarını biriktirdiniz.Diğer oyuncu şu mesajı gönderdi:} \\
\\ 
\\ 
\texttt{"(3,2) boş değilmiş, "┼" içeriyor. Benim tahtamda x=3 sütunundaki boş hücreler: (3,1) ve (3,6). Hedef düzeni şöyle değiştirelim: x=3 sütununda y=1'den y=7'ye: D,S,P,O,M,T,A. Yani M (3,5)'e, O (3,4)'e gelecek. M'ni (3,1)'e taşımak uygun mu?"} \\
\\ 
\\ 
\texttt{Bir sonraki komutun nedir?} \\
            }
        }
    }
    & \\ \\

    \theutterance \stepcounter{utterance}  
    & & & \multicolumn{4}{p{0.6\linewidth}}{
        \cellcolor[rgb]{0.9,0.9,0.9}{
            \makecell[{{p{\linewidth}}}]{
                \texttt{\tiny{[GM$\langle$P2]}}
                \texttt{SÖYLE: (3,1) boş değil, "T" harfi içeriyor. Benim tahtamda x=3 sütunundaki boş hücreler: (3,1) hariç hepsi dolu aslında. Tam boş hücreleri kontrol edeyim: (3,5) boş görünüyor. M (3,4)'ten (3,5)'e taşıyabilirim. Senin (3,5) boş mu? Eğer öyleyse, M'yi (3,5)'e taşımak uygun.} \\
            }
        }
    }
     \\ \\

    \theutterance \stepcounter{utterance}  
    & & & \multicolumn{2}{p{0.3\linewidth}}{
        \cellcolor[rgb]{0.9,0.9,0.9}{
            \makecell[{{p{\linewidth}}}]{
                \texttt{\tiny{[GM$|$GM]}}
                \texttt{SÖYLE: (3,1) boş değil, "T" harfi içeriyor. Benim tahtamda x=3 sütunundaki boş hücreler: (3,1) hariç hepsi dolu aslında. Tam boş hücreleri kontrol edeyim: (3,5) boş görünüyor. M (3,4)'ten (3,5)'e taşıyabilirim. Senin (3,5) boş mu? Eğer öyleyse, M'yi (3,5)'e taşımak uygun.} \\
            }
        }
    }
    & & \\ \\

    \theutterance \stepcounter{utterance}  
    & & \multicolumn{4}{p{0.6\linewidth}}{
        \cellcolor[rgb]{0.9,0.9,0.9}{
            \makecell[{{p{\linewidth}}}]{
                \texttt{\tiny{[P1$\langle$GM]}}
                \texttt{Mesajın diğer oyuncuya iletildi.} \\
\\ 
\\ 
\texttt{Şu anda en fazla 28 turdan 4. turu oynuyorsunuz.} \\
\texttt{Toplamda 16 cezanın 1 kadarını biriktirdiniz.Diğer oyuncu şu mesajı gönderdi:} \\
\\ 
\\ 
\texttt{"(3,1) boş değil, "T" harfi içeriyor. Benim tahtamda x=3 sütunundaki boş hücreler: (3,1) hariç hepsi dolu aslında. Tam boş hücreleri kontrol edeyim: (3,5) boş görünüyor. M (3,4)'ten (3,5)'e taşıyabilirim. Senin (3,5) boş mu? Eğer öyleyse, M'yi (3,5)'e taşımak uygun."} \\
\\ 
\\ 
\texttt{Bir sonraki komutun nedir?} \\
            }
        }
    }
    & \\ \\

    \theutterance \stepcounter{utterance}  
    & \multicolumn{4}{p{0.6\linewidth}}{
        \cellcolor[rgb]{0.9,0.9,0.9}{
            \makecell[{{p{\linewidth}}}]{
                \texttt{\tiny{[P1$\rangle$GM]}}
                \texttt{SÖYLE: (3,5) benim tahtamda boş değil, D var orada. Hedef olarak x=3 sütununda y=1'den y=7'ye şu düzeni öneriyorum: S,P,M,O,D,T,A. Yani: S(3,1), P(3,2), M(3,3), O(3,4), D(3,5), T(3,6), A(3,7). Senin S(3,3) ve A(3,7) zaten yerinde ama taşınması gerekiyor. Senin (3,1) boş mu?} \\
            }
        }
    }
    & & \\ \\

    \theutterance \stepcounter{utterance}  
    & & & \multicolumn{2}{p{0.3\linewidth}}{
        \cellcolor[rgb]{0.9,0.9,0.9}{
            \makecell[{{p{\linewidth}}}]{
                \texttt{\tiny{[GM$|$GM]}}
                \texttt{SÖYLE: (3,5) benim tahtamda boş değil, D var orada. Hedef olarak x=3 sütununda y=1'den y=7'ye şu düzeni öneriyorum: S,P,M,O,D,T,A. Yani: S(3,1), P(3,2), M(3,3), O(3,4), D(3,5), T(3,6), A(3,7). Senin S(3,3) ve A(3,7) zaten yerinde ama taşınması gerekiyor. Senin (3,1) boş mu?} \\
            }
        }
    }
    & & \\ \\

    \theutterance \stepcounter{utterance}  
    & & \multicolumn{4}{p{0.6\linewidth}}{
        \cellcolor[rgb]{0.9,0.9,0.9}{
            \makecell[{{p{\linewidth}}}]{
                \texttt{\tiny{[GM$\rangle$P2]}}
                \texttt{Mesajın diğer oyuncuya iletildi.} \\
\\ 
\\ 
\texttt{Şu anda en fazla 28 turdan 5. turu oynuyorsunuz.} \\
\texttt{Toplamda 16 cezanın 1 kadarını biriktirdiniz.Diğer oyuncu şu mesajı gönderdi:} \\
\\ 
\\ 
\texttt{"(3,5) benim tahtamda boş değil, D var orada. Hedef olarak x=3 sütununda y=1'den y=7'ye şu düzeni öneriyorum: S,P,M,O,D,T,A. Yani: S(3,1), P(3,2), M(3,3), O(3,4), D(3,5), T(3,6), A(3,7). Senin S(3,3) ve A(3,7) zaten yerinde ama taşınması gerekiyor. Senin (3,1) boş mu?"} \\
\\ 
\\ 
\texttt{Bir sonraki komutun nedir?} \\
            }
        }
    }
    & \\ \\

    \theutterance \stepcounter{utterance}  
    & & & \multicolumn{4}{p{0.6\linewidth}}{
        \cellcolor[rgb]{0.9,0.9,0.9}{
            \makecell[{{p{\linewidth}}}]{
                \texttt{\tiny{[GM$\langle$P2]}}
                \texttt{SÖYLE: (3,1) bende "T" ile dolu, boş değil. Benim x=3'teki dolu hücreler: (3,1)=T, (3,3)=S, (3,4)=M, (3,7)=A. Boş hücreler: (3,2) hariç hepsi, yani (3,5) ve (3,6) boş olmalı. Ama önce tam boş hücreleri netleştirelim. Bana tüm tahtamdaki boş hücreleri saymak yerine, şunu sormak istiyorum: x=3 sütununda sendeki boş hücreler hangileri?} \\
            }
        }
    }
     \\ \\

    \theutterance \stepcounter{utterance}  
    & & & \multicolumn{2}{p{0.3\linewidth}}{
        \cellcolor[rgb]{0.9,0.9,0.9}{
            \makecell[{{p{\linewidth}}}]{
                \texttt{\tiny{[GM$|$GM]}}
                \texttt{SÖYLE: (3,1) bende "T" ile dolu, boş değil. Benim x=3'teki dolu hücreler: (3,1)=T, (3,3)=S, (3,4)=M, (3,7)=A. Boş hücreler: (3,2) hariç hepsi, yani (3,5) ve (3,6) boş olmalı. Ama önce tam boş hücreleri netleştirelim. Bana tüm tahtamdaki boş hücreleri saymak yerine, şunu sormak istiyorum: x=3 sütununda sendeki boş hücreler hangileri?} \\
            }
        }
    }
    & & \\ \\

    \theutterance \stepcounter{utterance}  
    & & \multicolumn{4}{p{0.6\linewidth}}{
        \cellcolor[rgb]{0.9,0.9,0.9}{
            \makecell[{{p{\linewidth}}}]{
                \texttt{\tiny{[P1$\langle$GM]}}
                \texttt{Mesajın diğer oyuncuya iletildi.} \\
\\ 
\\ 
\texttt{Şu anda en fazla 28 turdan 5. turu oynuyorsunuz.} \\
\texttt{Toplamda 16 cezanın 1 kadarını biriktirdiniz.Diğer oyuncu şu mesajı gönderdi:} \\
\\ 
\\ 
\texttt{"(3,1) bende "T" ile dolu, boş değil. Benim x=3'teki dolu hücreler: (3,1)=T, (3,3)=S, (3,4)=M, (3,7)=A. Boş hücreler: (3,2) hariç hepsi, yani (3,5) ve (3,6) boş olmalı. Ama önce tam boş hücreleri netleştirelim. Bana tüm tahtamdaki boş hücreleri saymak yerine, şunu sormak istiyorum: x=3 sütununda sendeki boş hücreler hangileri?"} \\
\\ 
\\ 
\texttt{Bir sonraki komutun nedir?} \\
            }
        }
    }
    & \\ \\

    \theutterance \stepcounter{utterance}  
    & \multicolumn{4}{p{0.6\linewidth}}{
        \cellcolor[rgb]{0.9,0.9,0.9}{
            \makecell[{{p{\linewidth}}}]{
                \texttt{\tiny{[P1$\rangle$GM]}}
                \texttt{SÖYLE: Benim tahtamda x=3 sütununda boş hücreler: (3,1) ve (3,6). Diğerleri dolu: S(3,3), A(3,4), T(3,7). D(5,1) ve (6,1) de dolu. Hedef olarak x=3 sütununu tamamen kullanmak zor görünüşüyor. Bunun yerine x=6 sütununu deneyelim mi? Benim x=6'da sadece D(6,5) var, gerisi boş. Senin x=6'da ne var?} \\
            }
        }
    }
    & & \\ \\

    \theutterance \stepcounter{utterance}  
    & & & \multicolumn{2}{p{0.3\linewidth}}{
        \cellcolor[rgb]{0.9,0.9,0.9}{
            \makecell[{{p{\linewidth}}}]{
                \texttt{\tiny{[GM$|$GM]}}
                \texttt{SÖYLE: Benim tahtamda x=3 sütununda boş hücreler: (3,1) ve (3,6). Diğerleri dolu: S(3,3), A(3,4), T(3,7). D(5,1) ve (6,1) de dolu. Hedef olarak x=3 sütununu tamamen kullanmak zor görünüşüyor. Bunun yerine x=6 sütununu deneyelim mi? Benim x=6'da sadece D(6,5) var, gerisi boş. Senin x=6'da ne var?} \\
            }
        }
    }
    & & \\ \\

    \theutterance \stepcounter{utterance}  
    & & \multicolumn{4}{p{0.6\linewidth}}{
        \cellcolor[rgb]{0.9,0.9,0.9}{
            \makecell[{{p{\linewidth}}}]{
                \texttt{\tiny{[GM$\rangle$P2]}}
                \texttt{Mesajın diğer oyuncuya iletildi.} \\
\\ 
\\ 
\texttt{Şu anda en fazla 28 turdan 6. turu oynuyorsunuz.} \\
\texttt{Toplamda 16 cezanın 1 kadarını biriktirdiniz.Diğer oyuncu şu mesajı gönderdi:} \\
\\ 
\\ 
\texttt{"Benim tahtamda x=3 sütununda boş hücreler: (3,1) ve (3,6). Diğerleri dolu: S(3,3), A(3,4), T(3,7). D(5,1) ve (6,1) de dolu. Hedef olarak x=3 sütununu tamamen kullanmak zor görünüşüyor. Bunun yerine x=6 sütununu deneyelim mi? Benim x=6'da sadece D(6,5) var, gerisi boş. Senin x=6'da ne var?"} \\
\\ 
\\ 
\texttt{Bir sonraki komutun nedir?} \\
            }
        }
    }
    & \\ \\

    \theutterance \stepcounter{utterance}  
    & & & \multicolumn{4}{p{0.6\linewidth}}{
        \cellcolor[rgb]{0.9,0.9,0.9}{
            \makecell[{{p{\linewidth}}}]{
                \texttt{\tiny{[GM$\langle$P2]}}
                \texttt{SÖYLE: Benim x=6'da P(6,3) ve D(6,5) var. Diğer hücreler boş olmalı. x=6 hedefi iyi olabilir! Şöyle bir düzen öneriyorum: x=6'da y=1'den y=7'ye: D,S,P,O,M,T,A. Yani D(6,1), S(6,2), P(6,3), O(6,4), M(6,5), T(6,6), A(6,7). Benim P(6,3) ve D(6,5) zaten x=6'da. Senin D(6,5) de x=6'da! Ama benim D'yi (6,5)'ten (6,1)'e taşımam lazım, (6,1) bende boş mu kontrol edeyim. Senin (6,1), (6,2), (6,3), (6,4), (6,6), (6,7) boş mu?} \\
            }
        }
    }
     \\ \\

    \theutterance \stepcounter{utterance}  
    & & & \multicolumn{2}{p{0.3\linewidth}}{
        \cellcolor[rgb]{0.9,0.9,0.9}{
            \makecell[{{p{\linewidth}}}]{
                \texttt{\tiny{[GM$|$GM]}}
                \texttt{SÖYLE: Benim x=6'da P(6,3) ve D(6,5) var. Diğer hücreler boş olmalı. x=6 hedefi iyi olabilir! Şöyle bir düzen öneriyorum: x=6'da y=1'den y=7'ye: D,S,P,O,M,T,A. Yani D(6,1), S(6,2), P(6,3), O(6,4), M(6,5), T(6,6), A(6,7). Benim P(6,3) ve D(6,5) zaten x=6'da. Senin D(6,5) de x=6'da! Ama benim D'yi (6,5)'ten (6,1)'e taşımam lazım, (6,1) bende boş mu kontrol edeyim. Senin (6,1), (6,2), (6,3), (6,4), (6,6), (6,7) boş mu?} \\
            }
        }
    }
    & & \\ \\

    \theutterance \stepcounter{utterance}  
    & & \multicolumn{4}{p{0.6\linewidth}}{
        \cellcolor[rgb]{0.9,0.9,0.9}{
            \makecell[{{p{\linewidth}}}]{
                \texttt{\tiny{[P1$\langle$GM]}}
                \texttt{Mesajın diğer oyuncuya iletildi.} \\
\\ 
\\ 
\texttt{Şu anda en fazla 28 turdan 6. turu oynuyorsunuz.} \\
\texttt{Toplamda 16 cezanın 1 kadarını biriktirdiniz.Diğer oyuncu şu mesajı gönderdi:} \\
\\ 
\\ 
\texttt{"Benim x=6'da P(6,3) ve D(6,5) var. Diğer hücreler boş olmalı. x=6 hedefi iyi olabilir! Şöyle bir düzen öneriyorum: x=6'da y=1'den y=7'ye: D,S,P,O,M,T,A. Yani D(6,1), S(6,2), P(6,3), O(6,4), M(6,5), T(6,6), A(6,7). Benim P(6,3) ve D(6,5) zaten x=6'da. Senin D(6,5) de x=6'da! Ama benim D'yi (6,5)'ten (6,1)'e taşımam lazım, (6,1) bende boş mu kontrol edeyim. Senin (6,1), (6,2), (6,3), (6,4), (6,6), (6,7) boş mu?"} \\
\\ 
\\ 
\texttt{Bir sonraki komutun nedir?} \\
            }
        }
    }
    & \\ \\

    \theutterance \stepcounter{utterance}  
    & \multicolumn{4}{p{0.6\linewidth}}{
        \cellcolor[rgb]{0.9,0.9,0.9}{
            \makecell[{{p{\linewidth}}}]{
                \texttt{\tiny{[P1$\rangle$GM]}}
                \texttt{SÖYLE: Benim tahtamda x=6 sütunundaki boş hücreler: (6,1), (6,2), (6,3), (6,4), (6,7). (6,6) boş değil. Hedef düzen x=6'da: D(6,1), S(6,2), P(6,3), O(6,4), M(6,5), T(6,7), A(6,?) olamaz çünkü (6,6) dolu. Alternatif: x=6'da D(6,1), S(6,2), P(6,3), O(6,4), M(6,5), A(6,7) ve T'yi başka yere koyalım. Senin (6,6) boş mu?} \\
            }
        }
    }
    & & \\ \\

    \theutterance \stepcounter{utterance}  
    & & & \multicolumn{2}{p{0.3\linewidth}}{
        \cellcolor[rgb]{0.9,0.9,0.9}{
            \makecell[{{p{\linewidth}}}]{
                \texttt{\tiny{[GM$|$GM]}}
                \texttt{SÖYLE: Benim tahtamda x=6 sütunundaki boş hücreler: (6,1), (6,2), (6,3), (6,4), (6,7). (6,6) boş değil. Hedef düzen x=6'da: D(6,1), S(6,2), P(6,3), O(6,4), M(6,5), T(6,7), A(6,?) olamaz çünkü (6,6) dolu. Alternatif: x=6'da D(6,1), S(6,2), P(6,3), O(6,4), M(6,5), A(6,7) ve T'yi başka yere koyalım. Senin (6,6) boş mu?} \\
            }
        }
    }
    & & \\ \\

    \theutterance \stepcounter{utterance}  
    & & \multicolumn{4}{p{0.6\linewidth}}{
        \cellcolor[rgb]{0.9,0.9,0.9}{
            \makecell[{{p{\linewidth}}}]{
                \texttt{\tiny{[GM$\rangle$P2]}}
                \texttt{Mesajın diğer oyuncuya iletildi.} \\
\\ 
\\ 
\texttt{Şu anda en fazla 28 turdan 7. turu oynuyorsunuz.} \\
\texttt{Toplamda 16 cezanın 1 kadarını biriktirdiniz.Diğer oyuncu şu mesajı gönderdi:} \\
\\ 
\\ 
\texttt{"Benim tahtamda x=6 sütunundaki boş hücreler: (6,1), (6,2), (6,3), (6,4), (6,7). (6,6) boş değil. Hedef düzen x=6'da: D(6,1), S(6,2), P(6,3), O(6,4), M(6,5), T(6,7), A(6,?) olamaz çünkü (6,6) dolu. Alternatif: x=6'da D(6,1), S(6,2), P(6,3), O(6,4), M(6,5), A(6,7) ve T'yi başka yere koyalım. Senin (6,6) boş mu?"} \\
\\ 
\\ 
\texttt{Bir sonraki komutun nedir?} \\
            }
        }
    }
    & \\ \\

    \theutterance \stepcounter{utterance}  
    & & & \multicolumn{4}{p{0.6\linewidth}}{
        \cellcolor[rgb]{0.9,0.9,0.9}{
            \makecell[{{p{\linewidth}}}]{
                \texttt{\tiny{[GM$\langle$P2]}}
                \texttt{SÖYLE: Benim (6,6) boş değil, "┬" içeriyor. Yani T'yi (6,6)'ya koyamam. Benim x=6'da dolu hücreler: (6,3)=P, (6,5)=D, (6,6)=┬. Boş hücreler: (6,1), (6,2), (6,4), (6,7). Şu düzeni öneriyorum: x=6'da D(6,1), S(6,2), P(6,3), O(6,4), M(6,5), A(6,7). T'yi de x=6'ya değil, başka boş yere koyalım—örneğin senin ve benim boş olan (4,3)'e T(4,3)? Senin (4,3) boş mu?} \\
            }
        }
    }
     \\ \\

    \theutterance \stepcounter{utterance}  
    & & & \multicolumn{2}{p{0.3\linewidth}}{
        \cellcolor[rgb]{0.9,0.9,0.9}{
            \makecell[{{p{\linewidth}}}]{
                \texttt{\tiny{[GM$|$GM]}}
                \texttt{SÖYLE: Benim (6,6) boş değil, "┬" içeriyor. Yani T'yi (6,6)'ya koyamam. Benim x=6'da dolu hücreler: (6,3)=P, (6,5)=D, (6,6)=┬. Boş hücreler: (6,1), (6,2), (6,4), (6,7). Şu düzeni öneriyorum: x=6'da D(6,1), S(6,2), P(6,3), O(6,4), M(6,5), A(6,7). T'yi de x=6'ya değil, başka boş yere koyalım—örneğin senin ve benim boş olan (4,3)'e T(4,3)? Senin (4,3) boş mu?} \\
            }
        }
    }
    & & \\ \\

    \theutterance \stepcounter{utterance}  
    & & \multicolumn{4}{p{0.6\linewidth}}{
        \cellcolor[rgb]{0.9,0.9,0.9}{
            \makecell[{{p{\linewidth}}}]{
                \texttt{\tiny{[P1$\langle$GM]}}
                \texttt{Mesajın diğer oyuncuya iletildi.} \\
\\ 
\\ 
\texttt{Şu anda en fazla 28 turdan 7. turu oynuyorsunuz.} \\
\texttt{Toplamda 16 cezanın 1 kadarını biriktirdiniz.Diğer oyuncu şu mesajı gönderdi:} \\
\\ 
\\ 
\texttt{"Benim (6,6) boş değil, "┬" içeriyor. Yani T'yi (6,6)'ya koyamam. Benim x=6'da dolu hücreler: (6,3)=P, (6,5)=D, (6,6)=┬. Boş hücreler: (6,1), (6,2), (6,4), (6,7). Şu düzeni öneriyorum: x=6'da D(6,1), S(6,2), P(6,3), O(6,4), M(6,5), A(6,7). T'yi de x=6'ya değil, başka boş yere koyalım—örneğin senin ve benim boş olan (4,3)'e T(4,3)? Senin (4,3) boş mu?"} \\
\\ 
\\ 
\texttt{Bir sonraki komutun nedir?} \\
            }
        }
    }
    & \\ \\

    \theutterance \stepcounter{utterance}  
    & \multicolumn{4}{p{0.6\linewidth}}{
        \cellcolor[rgb]{0.9,0.9,0.9}{
            \makecell[{{p{\linewidth}}}]{
                \texttt{\tiny{[P1$\rangle$GM]}}
                \texttt{SÖYLE: (4,3) benim tahtamda boş değil, "A" var. Hedef düzen: x=6'da D(6,1), S(6,2), P(6,3), O(6,4), M(6,5), A(6,7) ve T'yi (7,1)'e koyalım. Benim (7,1) boş. Senin (7,1) boş mu? Onaylarsan hamlelere başlayalım.} \\
            }
        }
    }
    & & \\ \\

    \theutterance \stepcounter{utterance}  
    & & & \multicolumn{2}{p{0.3\linewidth}}{
        \cellcolor[rgb]{0.9,0.9,0.9}{
            \makecell[{{p{\linewidth}}}]{
                \texttt{\tiny{[GM$|$GM]}}
                \texttt{SÖYLE: (4,3) benim tahtamda boş değil, "A" var. Hedef düzen: x=6'da D(6,1), S(6,2), P(6,3), O(6,4), M(6,5), A(6,7) ve T'yi (7,1)'e koyalım. Benim (7,1) boş. Senin (7,1) boş mu? Onaylarsan hamlelere başlayalım.} \\
            }
        }
    }
    & & \\ \\

    \theutterance \stepcounter{utterance}  
    & & \multicolumn{4}{p{0.6\linewidth}}{
        \cellcolor[rgb]{0.9,0.9,0.9}{
            \makecell[{{p{\linewidth}}}]{
                \texttt{\tiny{[GM$\rangle$P2]}}
                \texttt{Mesajın diğer oyuncuya iletildi.} \\
\\ 
\\ 
\texttt{Şu anda en fazla 28 turdan 8. turu oynuyorsunuz.} \\
\texttt{Toplamda 16 cezanın 1 kadarını biriktirdiniz.Diğer oyuncu şu mesajı gönderdi:} \\
\\ 
\\ 
\texttt{"(4,3) benim tahtamda boş değil, "A" var. Hedef düzen: x=6'da D(6,1), S(6,2), P(6,3), O(6,4), M(6,5), A(6,7) ve T'yi (7,1)'e koyalım. Benim (7,1) boş. Senin (7,1) boş mu? Onaylarsan hamlelere başlayalım."} \\
\\ 
\\ 
\texttt{Bir sonraki komutun nedir?} \\
            }
        }
    }
    & \\ \\

    \theutterance \stepcounter{utterance}  
    & & & \multicolumn{4}{p{0.6\linewidth}}{
        \cellcolor[rgb]{0.9,0.9,0.9}{
            \makecell[{{p{\linewidth}}}]{
                \texttt{\tiny{[GM$\langle$P2]}}
                \texttt{SÖYLE: (7,1) bende boş! Hedef düzen onaylanmıştır:} \\
\texttt{{-} D(6,1), S(6,2), P(6,3), O(6,4), M(6,5), A(6,7), T(7,1)} \\
\\ 
\texttt{Şimdi hamlelerimi planlıyorum. İlk olarak D'yi (6,5)'ten (6,1)'e taşımaya çalışıyorum, ama yolu kontrol etmem lazım—tek hamlede taşıyabilir miyim? Kurallarda sadece hedefin boş olması gerekiyor, yol önemli değil sanırım. İlk hamlemi yapıyorum:} \\
\\ 
\texttt{HAMLE: D, (6,1)} \\
            }
        }
    }
     \\ \\

    \theutterance \stepcounter{utterance}  
    & & & \multicolumn{2}{p{0.3\linewidth}}{
        \cellcolor[rgb]{0.9,0.9,0.9}{
            \makecell[{{p{\linewidth}}}]{
                \texttt{\tiny{[GM$|$GM]}}
                \texttt{SÖYLE: (7,1) bende boş! Hedef düzen onaylanmıştır:} \\
\texttt{{-} D(6,1), S(6,2), P(6,3), O(6,4), M(6,5), A(6,7), T(7,1)} \\
\\ 
\texttt{Şimdi hamlelerimi planlıyorum. İlk olarak D'yi (6,5)'ten (6,1)'e taşımaya çalışıyorum, ama yolu kontrol etmem lazım—tek hamlede taşıyabilir miyim? Kurallarda sadece hedefin boş olması gerekiyor, yol önemli değil sanırım. İlk hamlemi yapıyorum:} \\
\\ 
\texttt{HAMLE: D, (6,1)} \\
            }
        }
    }
    & & \\ \\

    \theutterance \stepcounter{utterance}  
    & & & \multicolumn{2}{p{0.3\linewidth}}{
        \cellcolor[rgb]{0.9,0.9,0.9}{
            \makecell[{{p{\linewidth}}}]{
                \texttt{\tiny{[GM$|$GM]}}
                \texttt{Invalid response format: SÖYLE: (7,1) bende boş! Hedef düzen onaylanmıştır:} \\
\texttt{{-} D(6,1), S(6,2), P(6,3), O(6,4), M(6,5), A(6,7), T(7,1)} \\
\\ 
\texttt{Şimdi hamlelerimi planlıyorum. İlk olarak D'yi (6,5)'ten (6,1)'e taşımaya çalışıyorum, ama yolu kontrol etmem lazım—tek hamlede taşıyabilir miyim? Kurallarda sadece hedefin boş olması gerekiyor, yol önemli değil sanırım. İlk hamlemi yapıyorum:} \\
\\ 
\texttt{HAMLE: D, (6,1)} \\
            }
        }
    }
    & & \\ \\

    \theutterance \stepcounter{utterance}  
    & & \multicolumn{4}{p{0.6\linewidth}}{
        \cellcolor[rgb]{0.9,0.9,0.9}{
            \makecell[{{p{\linewidth}}}]{
                \texttt{\tiny{[GM$\rangle$P2]}}
                \texttt{Ceza: Mesajında birden fazla komut var!} \\
\\ 
\texttt{Yanıtının yalnızca SÖYLE: <MESAJ> ya da HAMLE: <NESNE>, (<X>, <Y>) komutlarından birini içerdiğinden emin ol; bunun dışında hiçbir şey yazma!} \\
\\ 
\texttt{Toplamda 16 cezanın 2 kadarını biriktirdiniz.Lütfen tekrar deneyin!} \\
            }
        }
    }
    & \\ \\

    \theutterance \stepcounter{utterance}  
    & & & \multicolumn{4}{p{0.6\linewidth}}{
        \cellcolor[rgb]{0.9,0.9,0.9}{
            \makecell[{{p{\linewidth}}}]{
                \texttt{\tiny{[GM$\langle$P2]}}
                \texttt{HAMLE: D, (6,1)} \\
            }
        }
    }
     \\ \\

    \theutterance \stepcounter{utterance}  
    & & & \multicolumn{2}{p{0.3\linewidth}}{
        \cellcolor[rgb]{0.9,0.9,0.9}{
            \makecell[{{p{\linewidth}}}]{
                \texttt{\tiny{[GM$|$GM]}}
                \texttt{HAMLE: D, (6,1)} \\
            }
        }
    }
    & & \\ \\

    \theutterance \stepcounter{utterance}  
    & & & \multicolumn{2}{p{0.3\linewidth}}{
        \cellcolor[rgb]{0.9,0.9,0.9}{
            \makecell[{{p{\linewidth}}}]{
                \texttt{\tiny{[GM$|$GM]}}
                \texttt{Ceza: (6,1) boş değil; "│" içeriyor.} \\
\\ 
\texttt{Toplamda 16 cezanın 3 kadarını biriktirdiniz.Lütfen tekrar deneyin!} \\
            }
        }
    }
    & & \\ \\

    \theutterance \stepcounter{utterance}  
    & & \multicolumn{4}{p{0.6\linewidth}}{
        \cellcolor[rgb]{0.9,0.9,0.9}{
            \makecell[{{p{\linewidth}}}]{
                \texttt{\tiny{[GM$\rangle$P2]}}
                \texttt{Ceza: (6,1) boş değil; "│" içeriyor.} \\
\\ 
\texttt{Toplamda 16 cezanın 3 kadarını biriktirdiniz.Lütfen tekrar deneyin!} \\
            }
        }
    }
    & \\ \\

    \theutterance \stepcounter{utterance}  
    & & & \multicolumn{4}{p{0.6\linewidth}}{
        \cellcolor[rgb]{0.9,0.9,0.9}{
            \makecell[{{p{\linewidth}}}]{
                \texttt{\tiny{[GM$\langle$P2]}}
                \texttt{SÖYLE: (6,1) bende "│" ile dolu, boş değil. Benim x=6'daki boş hücreleri tekrar kontrol etmem lazım. Tahtama bakınca x=6'da boş hücreler: (6,2) ve (6,4) ve (6,7) olabilir. (6,1) dolu. Hedef düzeni güncellememiz gerek. Senin x=6'daki boş hücrelerin hangileri, tekrar paylaşır mısın?} \\
            }
        }
    }
     \\ \\

    \theutterance \stepcounter{utterance}  
    & & & \multicolumn{2}{p{0.3\linewidth}}{
        \cellcolor[rgb]{0.9,0.9,0.9}{
            \makecell[{{p{\linewidth}}}]{
                \texttt{\tiny{[GM$|$GM]}}
                \texttt{SÖYLE: (6,1) bende "│" ile dolu, boş değil. Benim x=6'daki boş hücreleri tekrar kontrol etmem lazım. Tahtama bakınca x=6'da boş hücreler: (6,2) ve (6,4) ve (6,7) olabilir. (6,1) dolu. Hedef düzeni güncellememiz gerek. Senin x=6'daki boş hücrelerin hangileri, tekrar paylaşır mısın?} \\
            }
        }
    }
    & & \\ \\

    \theutterance \stepcounter{utterance}  
    & & \multicolumn{4}{p{0.6\linewidth}}{
        \cellcolor[rgb]{0.9,0.9,0.9}{
            \makecell[{{p{\linewidth}}}]{
                \texttt{\tiny{[P1$\langle$GM]}}
                \texttt{Mesajın diğer oyuncuya iletildi.} \\
\\ 
\\ 
\texttt{Şu anda en fazla 28 turdan 8. turu oynuyorsunuz.} \\
\texttt{Toplamda 16 cezanın 3 kadarını biriktirdiniz.Diğer oyuncu şu mesajı gönderdi:} \\
\\ 
\\ 
\texttt{"(6,1) bende "│" ile dolu, boş değil. Benim x=6'daki boş hücreleri tekrar kontrol etmem lazım. Tahtama bakınca x=6'da boş hücreler: (6,2) ve (6,4) ve (6,7) olabilir. (6,1) dolu. Hedef düzeni güncellememiz gerek. Senin x=6'daki boş hücrelerin hangileri, tekrar paylaşır mısın?"} \\
\\ 
\\ 
\texttt{Bir sonraki komutun nedir?} \\
            }
        }
    }
    & \\ \\

    \theutterance \stepcounter{utterance}  
    & \multicolumn{4}{p{0.6\linewidth}}{
        \cellcolor[rgb]{0.9,0.9,0.9}{
            \makecell[{{p{\linewidth}}}]{
                \texttt{\tiny{[P1$\rangle$GM]}}
                \texttt{SÖYLE: Benim x=6 boş hücreler: (6,1), (6,2), (6,3), (6,4), (6,7). (6,5)=D ve (6,6) dolu. Sadece 3 boş hücremiz ortak: (6,2), (6,4), (6,7). Yetersiz. Başka bir yaklaşım deneyelim: y=1 satırını hedefleyelim. Benim y=1'de P(2,1) ve M(7,1) dolu, diğerleri boş: (1,1),(3,1),(4,1),(5,1),(6,1). Senin y=1'de hangi hücreler boş?} \\
            }
        }
    }
    & & \\ \\

    \theutterance \stepcounter{utterance}  
    & & & \multicolumn{2}{p{0.3\linewidth}}{
        \cellcolor[rgb]{0.9,0.9,0.9}{
            \makecell[{{p{\linewidth}}}]{
                \texttt{\tiny{[GM$|$GM]}}
                \texttt{SÖYLE: Benim x=6 boş hücreler: (6,1), (6,2), (6,3), (6,4), (6,7). (6,5)=D ve (6,6) dolu. Sadece 3 boş hücremiz ortak: (6,2), (6,4), (6,7). Yetersiz. Başka bir yaklaşım deneyelim: y=1 satırını hedefleyelim. Benim y=1'de P(2,1) ve M(7,1) dolu, diğerleri boş: (1,1),(3,1),(4,1),(5,1),(6,1). Senin y=1'de hangi hücreler boş?} \\
            }
        }
    }
    & & \\ \\

    \theutterance \stepcounter{utterance}  
    & & \multicolumn{4}{p{0.6\linewidth}}{
        \cellcolor[rgb]{0.9,0.9,0.9}{
            \makecell[{{p{\linewidth}}}]{
                \texttt{\tiny{[GM$\rangle$P2]}}
                \texttt{Mesajın diğer oyuncuya iletildi.} \\
\\ 
\\ 
\texttt{Şu anda en fazla 28 turdan 9. turu oynuyorsunuz.} \\
\texttt{Toplamda 16 cezanın 3 kadarını biriktirdiniz.Diğer oyuncu şu mesajı gönderdi:} \\
\\ 
\\ 
\texttt{"Benim x=6 boş hücreler: (6,1), (6,2), (6,3), (6,4), (6,7). (6,5)=D ve (6,6) dolu. Sadece 3 boş hücremiz ortak: (6,2), (6,4), (6,7). Yetersiz. Başka bir yaklaşım deneyelim: y=1 satırını hedefleyelim. Benim y=1'de P(2,1) ve M(7,1) dolu, diğerleri boş: (1,1),(3,1),(4,1),(5,1),(6,1). Senin y=1'de hangi hücreler boş?"} \\
\\ 
\\ 
\texttt{Bir sonraki komutun nedir?} \\
            }
        }
    }
    & \\ \\

    \theutterance \stepcounter{utterance}  
    & & & \multicolumn{4}{p{0.6\linewidth}}{
        \cellcolor[rgb]{0.9,0.9,0.9}{
            \makecell[{{p{\linewidth}}}]{
                \texttt{\tiny{[GM$\langle$P2]}}
                \texttt{SÖYLE: Benim y=1 satırında (1,1) boş, (2,1) boş, (3,1)=T, (4,1) boş, (5,1) boş, (6,1) dolu "│", (7,1) boş. Yani boş hücreler: (1,1), (2,1), (4,1), (5,1), (7,1). Ortak boş hücreler y=1'de: (1,1), (2,1), (4,1), (5,1). Sadece 4 tane, 7 nesne için yetersiz. y=2 satırını deneyelim mi? Benim y=2'de hangi hücreler boş, kontrol edeyim—önce senin y=2'deki boş hücreler neler?} \\
            }
        }
    }
     \\ \\

    \theutterance \stepcounter{utterance}  
    & & & \multicolumn{2}{p{0.3\linewidth}}{
        \cellcolor[rgb]{0.9,0.9,0.9}{
            \makecell[{{p{\linewidth}}}]{
                \texttt{\tiny{[GM$|$GM]}}
                \texttt{SÖYLE: Benim y=1 satırında (1,1) boş, (2,1) boş, (3,1)=T, (4,1) boş, (5,1) boş, (6,1) dolu "│", (7,1) boş. Yani boş hücreler: (1,1), (2,1), (4,1), (5,1), (7,1). Ortak boş hücreler y=1'de: (1,1), (2,1), (4,1), (5,1). Sadece 4 tane, 7 nesne için yetersiz. y=2 satırını deneyelim mi? Benim y=2'de hangi hücreler boş, kontrol edeyim—önce senin y=2'deki boş hücreler neler?} \\
            }
        }
    }
    & & \\ \\

    \theutterance \stepcounter{utterance}  
    & & \multicolumn{4}{p{0.6\linewidth}}{
        \cellcolor[rgb]{0.9,0.9,0.9}{
            \makecell[{{p{\linewidth}}}]{
                \texttt{\tiny{[P1$\langle$GM]}}
                \texttt{Mesajın diğer oyuncuya iletildi.} \\
\\ 
\\ 
\texttt{Şu anda en fazla 28 turdan 9. turu oynuyorsunuz.} \\
\texttt{Toplamda 16 cezanın 3 kadarını biriktirdiniz.Diğer oyuncu şu mesajı gönderdi:} \\
\\ 
\\ 
\texttt{"Benim y=1 satırında (1,1) boş, (2,1) boş, (3,1)=T, (4,1) boş, (5,1) boş, (6,1) dolu "│", (7,1) boş. Yani boş hücreler: (1,1), (2,1), (4,1), (5,1), (7,1). Ortak boş hücreler y=1'de: (1,1), (2,1), (4,1), (5,1). Sadece 4 tane, 7 nesne için yetersiz. y=2 satırını deneyelim mi? Benim y=2'de hangi hücreler boş, kontrol edeyim—önce senin y=2'deki boş hücreler neler?"} \\
\\ 
\\ 
\texttt{Bir sonraki komutun nedir?} \\
            }
        }
    }
    & \\ \\

    \theutterance \stepcounter{utterance}  
    & \multicolumn{4}{p{0.6\linewidth}}{
        \cellcolor[rgb]{0.9,0.9,0.9}{
            \makecell[{{p{\linewidth}}}]{
                \texttt{\tiny{[P1$\rangle$GM]}}
                \texttt{SÖYLE: Benim y=2'de tüm hücreler dolu: (1,2)=P, (2,2)=M, (3,2)=┼, (4,2)=┴, (5,2)=◌, (6,2)=┼, (7,2)=─. Yetersiz. y=1'de 4 ortak boş hücremiz var: (1,1),(2,1),(4,1),(5,1). En yakın puan için bunları kullanalım. Hedef: 4 nesneyi y=1'de hizalayalım: S(1,1), P(2,1), O(4,1), T(5,1). Kalan 3 nesne D,M,A için de en yakın boş hücreleri bulalım. Senin (1,2),(2,2),(4,2) boş mü?} \\
            }
        }
    }
    & & \\ \\

    \theutterance \stepcounter{utterance}  
    & & & \multicolumn{2}{p{0.3\linewidth}}{
        \cellcolor[rgb]{0.9,0.9,0.9}{
            \makecell[{{p{\linewidth}}}]{
                \texttt{\tiny{[GM$|$GM]}}
                \texttt{SÖYLE: Benim y=2'de tüm hücreler dolu: (1,2)=P, (2,2)=M, (3,2)=┼, (4,2)=┴, (5,2)=◌, (6,2)=┼, (7,2)=─. Yetersiz. y=1'de 4 ortak boş hücremiz var: (1,1),(2,1),(4,1),(5,1). En yakın puan için bunları kullanalım. Hedef: 4 nesneyi y=1'de hizalayalım: S(1,1), P(2,1), O(4,1), T(5,1). Kalan 3 nesne D,M,A için de en yakın boş hücreleri bulalım. Senin (1,2),(2,2),(4,2) boş mü?} \\
            }
        }
    }
    & & \\ \\

    \theutterance \stepcounter{utterance}  
    & & \multicolumn{4}{p{0.6\linewidth}}{
        \cellcolor[rgb]{0.9,0.9,0.9}{
            \makecell[{{p{\linewidth}}}]{
                \texttt{\tiny{[GM$\rangle$P2]}}
                \texttt{Mesajın diğer oyuncuya iletildi.} \\
\\ 
\\ 
\texttt{Şu anda en fazla 28 turdan 10. turu oynuyorsunuz.} \\
\texttt{Toplamda 16 cezanın 3 kadarını biriktirdiniz.Diğer oyuncu şu mesajı gönderdi:} \\
\\ 
\\ 
\texttt{"Benim y=2'de tüm hücreler dolu: (1,2)=P, (2,2)=M, (3,2)=┼, (4,2)=┴, (5,2)=◌, (6,2)=┼, (7,2)=─. Yetersiz. y=1'de 4 ortak boş hücremiz var: (1,1),(2,1),(4,1),(5,1). En yakın puan için bunları kullanalım. Hedef: 4 nesneyi y=1'de hizalayalım: S(1,1), P(2,1), O(4,1), T(5,1). Kalan 3 nesne D,M,A için de en yakın boş hücreleri bulalım. Senin (1,2),(2,2),(4,2) boş mü?"} \\
\\ 
\\ 
\texttt{Bir sonraki komutun nedir?} \\
            }
        }
    }
    & \\ \\

    \theutterance \stepcounter{utterance}  
    & & & \multicolumn{4}{p{0.6\linewidth}}{
        \cellcolor[rgb]{0.9,0.9,0.9}{
            \makecell[{{p{\linewidth}}}]{
                \texttt{\tiny{[GM$\langle$P2]}}
                \texttt{SÖYLE: Benim (1,2) dolu "║", (2,2) dolu "──", (4,2) dolu "──". Yetersiz. y=1'de 4 ortak boş hücre var ama 7 nesne için yetersiz. Yeni yaklaşım: y=5 satırını deneyelim. Benim y=5'de D(6,5) var, diğerleri büyük olasılıkla boş. Senin y=5'deki dolu ve boş hücreler neler? Veya y=7 satırı—benim y=7'de A(3,7) var, diğerleri boş olabilir. Senin y=7 nasıl?} \\
            }
        }
    }
     \\ \\

    \theutterance \stepcounter{utterance}  
    & & & \multicolumn{2}{p{0.3\linewidth}}{
        \cellcolor[rgb]{0.9,0.9,0.9}{
            \makecell[{{p{\linewidth}}}]{
                \texttt{\tiny{[GM$|$GM]}}
                \texttt{SÖYLE: Benim (1,2) dolu "║", (2,2) dolu "──", (4,2) dolu "──". Yetersiz. y=1'de 4 ortak boş hücre var ama 7 nesne için yetersiz. Yeni yaklaşım: y=5 satırını deneyelim. Benim y=5'de D(6,5) var, diğerleri büyük olasılıkla boş. Senin y=5'deki dolu ve boş hücreler neler? Veya y=7 satırı—benim y=7'de A(3,7) var, diğerleri boş olabilir. Senin y=7 nasıl?} \\
            }
        }
    }
    & & \\ \\

    \theutterance \stepcounter{utterance}  
    & & \multicolumn{4}{p{0.6\linewidth}}{
        \cellcolor[rgb]{0.9,0.9,0.9}{
            \makecell[{{p{\linewidth}}}]{
                \texttt{\tiny{[P1$\langle$GM]}}
                \texttt{Mesajın diğer oyuncuya iletildi.} \\
\\ 
\\ 
\texttt{Şu anda en fazla 28 turdan 10. turu oynuyorsunuz.} \\
\texttt{Toplamda 16 cezanın 3 kadarını biriktirdiniz.Diğer oyuncu şu mesajı gönderdi:} \\
\\ 
\\ 
\texttt{"Benim (1,2) dolu "║", (2,2) dolu "──", (4,2) dolu "──". Yetersiz. y=1'de 4 ortak boş hücre var ama 7 nesne için yetersiz. Yeni yaklaşım: y=5 satırını deneyelim. Benim y=5'de D(6,5) var, diğerleri büyük olasılıkla boş. Senin y=5'deki dolu ve boş hücreler neler? Veya y=7 satırı—benim y=7'de A(3,7) var, diğerleri boş olabilir. Senin y=7 nasıl?"} \\
\\ 
\\ 
\texttt{Bir sonraki komutun nedir?} \\
            }
        }
    }
    & \\ \\

    \theutterance \stepcounter{utterance}  
    & \multicolumn{4}{p{0.6\linewidth}}{
        \cellcolor[rgb]{0.9,0.9,0.9}{
            \makecell[{{p{\linewidth}}}]{
                \texttt{\tiny{[P1$\rangle$GM]}}
                \texttt{SÖYLE: Benim y=5 satırım: (1,5)=D, (5,5)=◌, (6,5)=◌, diğerleri boş. y=7 satırım: (3,7)=T, (1,7)=T, (6,7)=│, diğerleri boş. y=5'te 5 ortak boş hücre olabilir. Senin y=5'teki tam dolu hücreler hangileri? Sadece (6,5)=D mi var?} \\
            }
        }
    }
    & & \\ \\

    \theutterance \stepcounter{utterance}  
    & & & \multicolumn{2}{p{0.3\linewidth}}{
        \cellcolor[rgb]{0.9,0.9,0.9}{
            \makecell[{{p{\linewidth}}}]{
                \texttt{\tiny{[GM$|$GM]}}
                \texttt{SÖYLE: Benim y=5 satırım: (1,5)=D, (5,5)=◌, (6,5)=◌, diğerleri boş. y=7 satırım: (3,7)=T, (1,7)=T, (6,7)=│, diğerleri boş. y=5'te 5 ortak boş hücre olabilir. Senin y=5'teki tam dolu hücreler hangileri? Sadece (6,5)=D mi var?} \\
            }
        }
    }
    & & \\ \\

    \theutterance \stepcounter{utterance}  
    & & \multicolumn{4}{p{0.6\linewidth}}{
        \cellcolor[rgb]{0.9,0.9,0.9}{
            \makecell[{{p{\linewidth}}}]{
                \texttt{\tiny{[GM$\rangle$P2]}}
                \texttt{Mesajın diğer oyuncuya iletildi.} \\
\\ 
\\ 
\texttt{Şu anda en fazla 28 turdan 11. turu oynuyorsunuz.} \\
\texttt{Toplamda 16 cezanın 3 kadarını biriktirdiniz.Diğer oyuncu şu mesajı gönderdi:} \\
\\ 
\\ 
\texttt{"Benim y=5 satırım: (1,5)=D, (5,5)=◌, (6,5)=◌, diğerleri boş. y=7 satırım: (3,7)=T, (1,7)=T, (6,7)=│, diğerleri boş. y=5'te 5 ortak boş hücre olabilir. Senin y=5'teki tam dolu hücreler hangileri? Sadece (6,5)=D mi var?"} \\
\\ 
\\ 
\texttt{Bir sonraki komutun nedir?} \\
            }
        }
    }
    & \\ \\

    \theutterance \stepcounter{utterance}  
    & & & \multicolumn{4}{p{0.6\linewidth}}{
        \cellcolor[rgb]{0.9,0.9,0.9}{
            \makecell[{{p{\linewidth}}}]{
                \texttt{\tiny{[GM$\langle$P2]}}
                \texttt{SÖYLE: Benim y=5 satırımda (6,5)=D var. Diğer hücreleri tek tek kontrol ediyorum: (1,5) boş, (2,5) boş, (3,5) boş, (4,5) boş, (5,5) boş, (7,5) boş. Yani sadece (6,5)=D dolu, gerisi boş! 6 ortak boş hücre var. Hedef düzen y=5'te: S(1,5), P(2,5), M(3,5), O(4,5), A(5,5), D(6,5). T'yi de (7,5)'e koyalım—senin (7,5) boş mu? Onaylarsan hamlelere başlarız.} \\
            }
        }
    }
     \\ \\

    \theutterance \stepcounter{utterance}  
    & & & \multicolumn{2}{p{0.3\linewidth}}{
        \cellcolor[rgb]{0.9,0.9,0.9}{
            \makecell[{{p{\linewidth}}}]{
                \texttt{\tiny{[GM$|$GM]}}
                \texttt{SÖYLE: Benim y=5 satırımda (6,5)=D var. Diğer hücreleri tek tek kontrol ediyorum: (1,5) boş, (2,5) boş, (3,5) boş, (4,5) boş, (5,5) boş, (7,5) boş. Yani sadece (6,5)=D dolu, gerisi boş! 6 ortak boş hücre var. Hedef düzen y=5'te: S(1,5), P(2,5), M(3,5), O(4,5), A(5,5), D(6,5). T'yi de (7,5)'e koyalım—senin (7,5) boş mu? Onaylarsan hamlelere başlarız.} \\
            }
        }
    }
    & & \\ \\

    \theutterance \stepcounter{utterance}  
    & & \multicolumn{4}{p{0.6\linewidth}}{
        \cellcolor[rgb]{0.9,0.9,0.9}{
            \makecell[{{p{\linewidth}}}]{
                \texttt{\tiny{[P1$\langle$GM]}}
                \texttt{Mesajın diğer oyuncuya iletildi.} \\
\\ 
\\ 
\texttt{Şu anda en fazla 28 turdan 11. turu oynuyorsunuz.} \\
\texttt{Toplamda 16 cezanın 3 kadarını biriktirdiniz.Diğer oyuncu şu mesajı gönderdi:} \\
\\ 
\\ 
\texttt{"Benim y=5 satırımda (6,5)=D var. Diğer hücreleri tek tek kontrol ediyorum: (1,5) boş, (2,5) boş, (3,5) boş, (4,5) boş, (5,5) boş, (7,5) boş. Yani sadece (6,5)=D dolu, gerisi boş! 6 ortak boş hücre var. Hedef düzen y=5'te: S(1,5), P(2,5), M(3,5), O(4,5), A(5,5), D(6,5). T'yi de (7,5)'e koyalım—senin (7,5) boş mu? Onaylarsan hamlelere başlarız."} \\
\\ 
\\ 
\texttt{Bir sonraki komutun nedir?} \\
            }
        }
    }
    & \\ \\

    \theutterance \stepcounter{utterance}  
    & \multicolumn{4}{p{0.6\linewidth}}{
        \cellcolor[rgb]{0.9,0.9,0.9}{
            \makecell[{{p{\linewidth}}}]{
                \texttt{\tiny{[P1$\rangle$GM]}}
                \texttt{SÖYLE: (7,5) benim tahtamda boş değil, "│" içeriyor. Ama (1,5) de bende dolu, D var orada. Hedef düzeni şöyle yapalım: P(2,5), S(3,5), M(4,5), O(5,5), A(7,5), T'yi (1,5) yerine başka yere koyalım. Senin (3,5) ve (4,5) kesin boş mu? T için (7,1) deneyelim—benim (7,1) boş. Senin (7,1) boş mu?} \\
            }
        }
    }
    & & \\ \\

    \theutterance \stepcounter{utterance}  
    & & & \multicolumn{2}{p{0.3\linewidth}}{
        \cellcolor[rgb]{0.9,0.9,0.9}{
            \makecell[{{p{\linewidth}}}]{
                \texttt{\tiny{[GM$|$GM]}}
                \texttt{SÖYLE: (7,5) benim tahtamda boş değil, "│" içeriyor. Ama (1,5) de bende dolu, D var orada. Hedef düzeni şöyle yapalım: P(2,5), S(3,5), M(4,5), O(5,5), A(7,5), T'yi (1,5) yerine başka yere koyalım. Senin (3,5) ve (4,5) kesin boş mu? T için (7,1) deneyelim—benim (7,1) boş. Senin (7,1) boş mu?} \\
            }
        }
    }
    & & \\ \\

    \theutterance \stepcounter{utterance}  
    & & \multicolumn{4}{p{0.6\linewidth}}{
        \cellcolor[rgb]{0.9,0.9,0.9}{
            \makecell[{{p{\linewidth}}}]{
                \texttt{\tiny{[GM$\rangle$P2]}}
                \texttt{Mesajın diğer oyuncuya iletildi.} \\
\\ 
\\ 
\texttt{Şu anda en fazla 28 turdan 12. turu oynuyorsunuz.} \\
\texttt{Toplamda 16 cezanın 3 kadarını biriktirdiniz.Diğer oyuncu şu mesajı gönderdi:} \\
\\ 
\\ 
\texttt{"(7,5) benim tahtamda boş değil, "│" içeriyor. Ama (1,5) de bende dolu, D var orada. Hedef düzeni şöyle yapalım: P(2,5), S(3,5), M(4,5), O(5,5), A(7,5), T'yi (1,5) yerine başka yere koyalım. Senin (3,5) ve (4,5) kesin boş mu? T için (7,1) deneyelim—benim (7,1) boş. Senin (7,1) boş mu?"} \\
\\ 
\\ 
\texttt{Bir sonraki komutun nedir?} \\
            }
        }
    }
    & \\ \\

    \theutterance \stepcounter{utterance}  
    & & & \multicolumn{4}{p{0.6\linewidth}}{
        \cellcolor[rgb]{0.9,0.9,0.9}{
            \makecell[{{p{\linewidth}}}]{
                \texttt{\tiny{[GM$\langle$P2]}}
                \texttt{SÖYLE: (7,1) bende boş, evet! Ama (7,5) bende de dolu "│" ile, T'yi oraya koyamam. (3,5) ve (4,5) bende boş. Şu hedef düzeni öneriyorum: P(2,5), S(3,5), M(4,5), O(5,5), A(7,5 yerine A(1,5) deneyelim—senin (1,5) dolu D, olmaz. A'yı nereye koyalım? Senin y=5'teki boş hücreler tam olarak hangileri? Benim y=5 boş hücreler: (1,5),(2,5),(3,5),(4,5),(5,5),(7,5). (6,5)=D dolu.} \\
            }
        }
    }
     \\ \\

    \theutterance \stepcounter{utterance}  
    & & & \multicolumn{2}{p{0.3\linewidth}}{
        \cellcolor[rgb]{0.9,0.9,0.9}{
            \makecell[{{p{\linewidth}}}]{
                \texttt{\tiny{[GM$|$GM]}}
                \texttt{SÖYLE: (7,1) bende boş, evet! Ama (7,5) bende de dolu "│" ile, T'yi oraya koyamam. (3,5) ve (4,5) bende boş. Şu hedef düzeni öneriyorum: P(2,5), S(3,5), M(4,5), O(5,5), A(7,5 yerine A(1,5) deneyelim—senin (1,5) dolu D, olmaz. A'yı nereye koyalım? Senin y=5'teki boş hücreler tam olarak hangileri? Benim y=5 boş hücreler: (1,5),(2,5),(3,5),(4,5),(5,5),(7,5). (6,5)=D dolu.} \\
            }
        }
    }
    & & \\ \\

    \theutterance \stepcounter{utterance}  
    & & \multicolumn{4}{p{0.6\linewidth}}{
        \cellcolor[rgb]{0.9,0.9,0.9}{
            \makecell[{{p{\linewidth}}}]{
                \texttt{\tiny{[P1$\langle$GM]}}
                \texttt{Mesajın diğer oyuncuya iletildi.} \\
\\ 
\\ 
\texttt{Şu anda en fazla 28 turdan 12. turu oynuyorsunuz.} \\
\texttt{Toplamda 16 cezanın 3 kadarını biriktirdiniz.Diğer oyuncu şu mesajı gönderdi:} \\
\\ 
\\ 
\texttt{"(7,1) bende boş, evet! Ama (7,5) bende de dolu "│" ile, T'yi oraya koyamam. (3,5) ve (4,5) bende boş. Şu hedef düzeni öneriyorum: P(2,5), S(3,5), M(4,5), O(5,5), A(7,5 yerine A(1,5) deneyelim—senin (1,5) dolu D, olmaz. A'yı nereye koyalım? Senin y=5'teki boş hücreler tam olarak hangileri? Benim y=5 boş hücreler: (1,5),(2,5),(3,5),(4,5),(5,5),(7,5). (6,5)=D dolu."} \\
\\ 
\\ 
\texttt{Bir sonraki komutun nedir?} \\
            }
        }
    }
    & \\ \\

    \theutterance \stepcounter{utterance}  
    & \multicolumn{4}{p{0.6\linewidth}}{
        \cellcolor[rgb]{0.9,0.9,0.9}{
            \makecell[{{p{\linewidth}}}]{
                \texttt{\tiny{[P1$\rangle$GM]}}
                \texttt{SÖYLE: Benim y=5 boş hücreler: (2,5), (3,5), (4,5), (7,5). (1,5)=D ve (5,5)=◌ ve (6,5)=D dolu. (5,5) boş değil "◌" içeriyor ama sembol mü nesne mu belirsiz. Ortak boş hücreler: (2,5), (3,5), (4,5), (7,5) — sadece 4 tane. Hedef: P(2,5), S(3,5), M(4,5), O(7,5). T'yi (7,1)'e koyalım. D(6,5) ve A zaten yerinde kalabilir. Onaylarsan hamlelere başlıyorum, P'yi (2,1)'den (2,5)'e taşıyacağım.} \\
            }
        }
    }
    & & \\ \\

    \theutterance \stepcounter{utterance}  
    & & & \multicolumn{2}{p{0.3\linewidth}}{
        \cellcolor[rgb]{0.9,0.9,0.9}{
            \makecell[{{p{\linewidth}}}]{
                \texttt{\tiny{[GM$|$GM]}}
                \texttt{SÖYLE: Benim y=5 boş hücreler: (2,5), (3,5), (4,5), (7,5). (1,5)=D ve (5,5)=◌ ve (6,5)=D dolu. (5,5) boş değil "◌" içeriyor ama sembol mü nesne mu belirsiz. Ortak boş hücreler: (2,5), (3,5), (4,5), (7,5) — sadece 4 tane. Hedef: P(2,5), S(3,5), M(4,5), O(7,5). T'yi (7,1)'e koyalım. D(6,5) ve A zaten yerinde kalabilir. Onaylarsan hamlelere başlıyorum, P'yi (2,1)'den (2,5)'e taşıyacağım.} \\
            }
        }
    }
    & & \\ \\

    \theutterance \stepcounter{utterance}  
    & & \multicolumn{4}{p{0.6\linewidth}}{
        \cellcolor[rgb]{0.9,0.9,0.9}{
            \makecell[{{p{\linewidth}}}]{
                \texttt{\tiny{[GM$\rangle$P2]}}
                \texttt{Mesajın diğer oyuncuya iletildi.} \\
\\ 
\\ 
\texttt{Şu anda en fazla 28 turdan 13. turu oynuyorsunuz.} \\
\texttt{Toplamda 16 cezanın 3 kadarını biriktirdiniz.Diğer oyuncu şu mesajı gönderdi:} \\
\\ 
\\ 
\texttt{"Benim y=5 boş hücreler: (2,5), (3,5), (4,5), (7,5). (1,5)=D ve (5,5)=◌ ve (6,5)=D dolu. (5,5) boş değil "◌" içeriyor ama sembol mü nesne mu belirsiz. Ortak boş hücreler: (2,5), (3,5), (4,5), (7,5) — sadece 4 tane. Hedef: P(2,5), S(3,5), M(4,5), O(7,5). T'yi (7,1)'e koyalım. D(6,5) ve A zaten yerinde kalabilir. Onaylarsan hamlelere başlıyorum, P'yi (2,1)'den (2,5)'e taşıyacağım."} \\
\\ 
\\ 
\texttt{Bir sonraki komutun nedir?} \\
            }
        }
    }
    & \\ \\

    \theutterance \stepcounter{utterance}  
    & & & \multicolumn{4}{p{0.6\linewidth}}{
        \cellcolor[rgb]{0.9,0.9,0.9}{
            \makecell[{{p{\linewidth}}}]{
                \texttt{\tiny{[GM$\langle$P2]}}
                \texttt{SÖYLE: Onaylıyorum! Hedef düzen:} \\
\texttt{{-} P(2,5), S(3,5), M(4,5), O(7,5), T(7,1)} \\
\texttt{{-} D(6,5) zaten yerinde, A(3,7) olduğu yerde kalsın.} \\
\\ 
\texttt{Ben başlıyorum, S'yi (3,3)'ten (3,5)'e taşıyorum.} \\
            }
        }
    }
     \\ \\

    \theutterance \stepcounter{utterance}  
    & & & \multicolumn{2}{p{0.3\linewidth}}{
        \cellcolor[rgb]{0.9,0.9,0.9}{
            \makecell[{{p{\linewidth}}}]{
                \texttt{\tiny{[GM$|$GM]}}
                \texttt{SÖYLE: Onaylıyorum! Hedef düzen:} \\
\texttt{{-} P(2,5), S(3,5), M(4,5), O(7,5), T(7,1)} \\
\texttt{{-} D(6,5) zaten yerinde, A(3,7) olduğu yerde kalsın.} \\
\\ 
\texttt{Ben başlıyorum, S'yi (3,3)'ten (3,5)'e taşıyorum.} \\
            }
        }
    }
    & & \\ \\

    \theutterance \stepcounter{utterance}  
    & & & \multicolumn{2}{p{0.3\linewidth}}{
        \cellcolor[rgb]{0.9,0.9,0.9}{
            \makecell[{{p{\linewidth}}}]{
                \texttt{\tiny{[GM$|$GM]}}
                \texttt{Invalid format: tail is not empty: '} \\
\texttt{{-} P(2,5), S(3,5), M(4,5), O(7,5), T(7,1)} \\
\texttt{{-} D(6,5) zaten yerinde, A(3,7) olduğu yerde kalsın.} \\
\\ 
\texttt{Ben başlıyorum, S'yi (3,3)'ten (3,5)'e taşıyorum.} \\
            }
        }
    }
    & & \\ \\

    \theutterance \stepcounter{utterance}  
    & & \multicolumn{4}{p{0.6\linewidth}}{
        \cellcolor[rgb]{0.9,0.9,0.9}{
            \makecell[{{p{\linewidth}}}]{
                \texttt{\tiny{[GM$\rangle$P2]}}
                \texttt{Ceza: Mesajında komuttan sonra hiçbir şey bulunmamalı!} \\
\\ 
\texttt{Yanıtının yalnızca SÖYLE: <MESAJ> ya da HAMLE: <NESNE>, (<X>, <Y>) komutlarından birini içerdiğinden emin ol; bunun dışında hiçbir şey yazma!} \\
\\ 
\texttt{Toplamda 16 cezanın 4 kadarını biriktirdiniz.Lütfen tekrar deneyin!} \\
            }
        }
    }
    & \\ \\

    \theutterance \stepcounter{utterance}  
    & & & \multicolumn{4}{p{0.6\linewidth}}{
        \cellcolor[rgb]{0.9,0.9,0.9}{
            \makecell[{{p{\linewidth}}}]{
                \texttt{\tiny{[GM$\langle$P2]}}
                \texttt{HAMLE: S, (3,5)} \\
            }
        }
    }
     \\ \\

    \theutterance \stepcounter{utterance}  
    & & & \multicolumn{2}{p{0.3\linewidth}}{
        \cellcolor[rgb]{0.9,0.9,0.9}{
            \makecell[{{p{\linewidth}}}]{
                \texttt{\tiny{[GM$|$GM]}}
                \texttt{HAMLE: S, (3,5)} \\
            }
        }
    }
    & & \\ \\

    \theutterance \stepcounter{utterance}  
    & & & \multicolumn{2}{p{0.3\linewidth}}{
        \cellcolor[rgb]{0.9,0.9,0.9}{
            \makecell[{{p{\linewidth}}}]{
                \texttt{\tiny{[GM$|$GM]}}
                \texttt{S moved to (3, 5)} \\
            }
        }
    }
    & & \\ \\

    \theutterance \stepcounter{utterance}  
    & & \multicolumn{4}{p{0.6\linewidth}}{
        \cellcolor[rgb]{0.9,0.9,0.9}{
            \makecell[{{p{\linewidth}}}]{
                \texttt{\tiny{[P1$\langle$GM]}}
                \texttt{Mesajın diğer oyuncuya iletildi.} \\
\\ 
\\ 
\texttt{Şu anda en fazla 28 turdan 13. turu oynuyorsunuz.} \\
\texttt{Toplamda 16 cezanın 4 kadarını biriktirdiniz.Diğer oyuncu kendi ızgarasında bir nesneyi hareket ettirdi.} \\
\\ 
\\ 
\texttt{Bir sonraki komutun nedir?} \\
            }
        }
    }
    & \\ \\

    \theutterance \stepcounter{utterance}  
    & \multicolumn{4}{p{0.6\linewidth}}{
        \cellcolor[rgb]{0.9,0.9,0.9}{
            \makecell[{{p{\linewidth}}}]{
                \texttt{\tiny{[P1$\rangle$GM]}}
                \texttt{HAMLE: P, (2, 5)} \\
            }
        }
    }
    & & \\ \\

    \theutterance \stepcounter{utterance}  
    & & & \multicolumn{2}{p{0.3\linewidth}}{
        \cellcolor[rgb]{0.9,0.9,0.9}{
            \makecell[{{p{\linewidth}}}]{
                \texttt{\tiny{[GM$|$GM]}}
                \texttt{HAMLE: P, (2, 5)} \\
            }
        }
    }
    & & \\ \\

    \theutterance \stepcounter{utterance}  
    & & & \multicolumn{2}{p{0.3\linewidth}}{
        \cellcolor[rgb]{0.9,0.9,0.9}{
            \makecell[{{p{\linewidth}}}]{
                \texttt{\tiny{[GM$|$GM]}}
                \texttt{P moved to (2, 5)} \\
            }
        }
    }
    & & \\ \\

    \theutterance \stepcounter{utterance}  
    & & \multicolumn{4}{p{0.6\linewidth}}{
        \cellcolor[rgb]{0.9,0.9,0.9}{
            \makecell[{{p{\linewidth}}}]{
                \texttt{\tiny{[GM$\rangle$P2]}}
                \texttt{S’ başarıyla (3,5) koordinatlarına taşındı. Güncellenmiş ızgaran şu şekilde:} \\
\texttt{```} \\
\texttt{ 1234567} \\
\texttt{╔══╤══╤═╗} \\
\texttt{║◌◌│◌◌│T║ 1} \\
\texttt{╟──┼──┴─╢ 2} \\
\texttt{║◌◌│◌◌◌P║ 3} \\
\texttt{╟──┘◌◌MO║ 4} \\
\texttt{║D◌S◌◌◌◌║ 5} \\
\texttt{║◌◌┌─┬──╢ 6} \\
\texttt{║◌◌│A│◌◌║ 7} \\
\texttt{╚══╧═╧══╝} \\
\texttt{```} \\
\texttt{Şu anda en fazla 28 turdan 14. turu oynuyorsunuz.} \\
\texttt{Toplamda 16 cezanın 4 kadarını biriktirdiniz.Diğer oyuncu kendi ızgarasında bir nesneyi hareket ettirdi.} \\
\\ 
\\ 
\texttt{Bir sonraki komutun nedir?} \\
            }
        }
    }
    & \\ \\

    \theutterance \stepcounter{utterance}  
    & & & \multicolumn{4}{p{0.6\linewidth}}{
        \cellcolor[rgb]{0.9,0.9,0.9}{
            \makecell[{{p{\linewidth}}}]{
                \texttt{\tiny{[GM$\langle$P2]}}
                \texttt{HAMLE: M, (4,5)} \\
            }
        }
    }
     \\ \\

    \theutterance \stepcounter{utterance}  
    & & & \multicolumn{2}{p{0.3\linewidth}}{
        \cellcolor[rgb]{0.9,0.9,0.9}{
            \makecell[{{p{\linewidth}}}]{
                \texttt{\tiny{[GM$|$GM]}}
                \texttt{HAMLE: M, (4,5)} \\
            }
        }
    }
    & & \\ \\

    \theutterance \stepcounter{utterance}  
    & & & \multicolumn{2}{p{0.3\linewidth}}{
        \cellcolor[rgb]{0.9,0.9,0.9}{
            \makecell[{{p{\linewidth}}}]{
                \texttt{\tiny{[GM$|$GM]}}
                \texttt{M moved to (4, 5)} \\
            }
        }
    }
    & & \\ \\

    \theutterance \stepcounter{utterance}  
    & & \multicolumn{4}{p{0.6\linewidth}}{
        \cellcolor[rgb]{0.9,0.9,0.9}{
            \makecell[{{p{\linewidth}}}]{
                \texttt{\tiny{[P1$\langle$GM]}}
                \texttt{P’ başarıyla (2,5) koordinatlarına taşındı. Güncellenmiş ızgaran şu şekilde:} \\
\texttt{```} \\
\texttt{ 1234567} \\
\texttt{╔══╤══╤═╗} \\
\texttt{║◌◌│◌M│◌║ 1} \\
\texttt{╟──┼──┴─╢ 2} \\
\texttt{║◌◌│◌S◌O║ 3} \\
\texttt{╟──┘◌A◌◌║ 4} \\
\texttt{║DP◌◌◌◌◌║ 5} \\
\texttt{║◌◌┌─┬──╢ 6} \\
\texttt{║T◌│◌│◌◌║ 7} \\
\texttt{╚══╧═╧══╝} \\
\texttt{```} \\
\texttt{Şu anda en fazla 28 turdan 14. turu oynuyorsunuz.} \\
\texttt{Toplamda 16 cezanın 4 kadarını biriktirdiniz.Diğer oyuncu kendi ızgarasında bir nesneyi hareket ettirdi.} \\
\\ 
\\ 
\texttt{Bir sonraki komutun nedir?} \\
            }
        }
    }
    & \\ \\

    \theutterance \stepcounter{utterance}  
    & \multicolumn{4}{p{0.6\linewidth}}{
        \cellcolor[rgb]{0.9,0.9,0.9}{
            \makecell[{{p{\linewidth}}}]{
                \texttt{\tiny{[P1$\rangle$GM]}}
                \texttt{SÖYLE: P'yi (2,5)'e başarıyla taşıdım. Şimdi S'yi (3,3)'ten (3,5)'e taşımaya hazırlanıyorum. Senin durum nedir?} \\
            }
        }
    }
    & & \\ \\

    \theutterance \stepcounter{utterance}  
    & & & \multicolumn{2}{p{0.3\linewidth}}{
        \cellcolor[rgb]{0.9,0.9,0.9}{
            \makecell[{{p{\linewidth}}}]{
                \texttt{\tiny{[GM$|$GM]}}
                \texttt{SÖYLE: P'yi (2,5)'e başarıyla taşıdım. Şimdi S'yi (3,3)'ten (3,5)'e taşımaya hazırlanıyorum. Senin durum nedir?} \\
            }
        }
    }
    & & \\ \\

    \theutterance \stepcounter{utterance}  
    & & \multicolumn{4}{p{0.6\linewidth}}{
        \cellcolor[rgb]{0.9,0.9,0.9}{
            \makecell[{{p{\linewidth}}}]{
                \texttt{\tiny{[GM$\rangle$P2]}}
                \texttt{M’ başarıyla (4,5) koordinatlarına taşındı. Güncellenmiş ızgaran şu şekilde:} \\
\texttt{```} \\
\texttt{ 1234567} \\
\texttt{╔══╤══╤═╗} \\
\texttt{║◌◌│◌◌│T║ 1} \\
\texttt{╟──┼──┴─╢ 2} \\
\texttt{║◌◌│◌◌◌P║ 3} \\
\texttt{╟──┘◌◌◌O║ 4} \\
\texttt{║D◌SM◌◌◌║ 5} \\
\texttt{║◌◌┌─┬──╢ 6} \\
\texttt{║◌◌│A│◌◌║ 7} \\
\texttt{╚══╧═╧══╝} \\
\texttt{```} \\
\texttt{Şu anda en fazla 28 turdan 15. turu oynuyorsunuz.} \\
\texttt{Toplamda 16 cezanın 4 kadarını biriktirdiniz.Diğer oyuncu şu mesajı gönderdi:} \\
\\ 
\\ 
\texttt{"P'yi (2,5)'e başarıyla taşıdım. Şimdi S'yi (3,3)'ten (3,5)'e taşımaya hazırlanıyorum. Senin durum nedir?"} \\
\\ 
\\ 
\texttt{Bir sonraki komutun nedir?} \\
            }
        }
    }
    & \\ \\

    \theutterance \stepcounter{utterance}  
    & & & \multicolumn{4}{p{0.6\linewidth}}{
        \cellcolor[rgb]{0.9,0.9,0.9}{
            \makecell[{{p{\linewidth}}}]{
                \texttt{\tiny{[GM$\langle$P2]}}
                \texttt{SÖYLE: S ve M zaten yerinde! S(3,5) ve M(4,5) taşındı. Şimdi P'yi (6,3)'ten (2,5)'e taşıyorum.} \\
            }
        }
    }
     \\ \\

    \theutterance \stepcounter{utterance}  
    & & & \multicolumn{2}{p{0.3\linewidth}}{
        \cellcolor[rgb]{0.9,0.9,0.9}{
            \makecell[{{p{\linewidth}}}]{
                \texttt{\tiny{[GM$|$GM]}}
                \texttt{SÖYLE: S ve M zaten yerinde! S(3,5) ve M(4,5) taşındı. Şimdi P'yi (6,3)'ten (2,5)'e taşıyorum.} \\
            }
        }
    }
    & & \\ \\

    \theutterance \stepcounter{utterance}  
    & & \multicolumn{4}{p{0.6\linewidth}}{
        \cellcolor[rgb]{0.9,0.9,0.9}{
            \makecell[{{p{\linewidth}}}]{
                \texttt{\tiny{[P1$\langle$GM]}}
                \texttt{Mesajın diğer oyuncuya iletildi.} \\
\\ 
\\ 
\texttt{Şu anda en fazla 28 turdan 15. turu oynuyorsunuz.} \\
\texttt{Toplamda 16 cezanın 4 kadarını biriktirdiniz.Diğer oyuncu şu mesajı gönderdi:} \\
\\ 
\\ 
\texttt{"S ve M zaten yerinde! S(3,5) ve M(4,5) taşındı. Şimdi P'yi (6,3)'ten (2,5)'e taşıyorum."} \\
\\ 
\\ 
\texttt{Bir sonraki komutun nedir?} \\
            }
        }
    }
    & \\ \\

    \theutterance \stepcounter{utterance}  
    & \multicolumn{4}{p{0.6\linewidth}}{
        \cellcolor[rgb]{0.9,0.9,0.9}{
            \makecell[{{p{\linewidth}}}]{
                \texttt{\tiny{[P1$\rangle$GM]}}
                \texttt{HAMLE: S, (3, 5)} \\
            }
        }
    }
    & & \\ \\

    \theutterance \stepcounter{utterance}  
    & & & \multicolumn{2}{p{0.3\linewidth}}{
        \cellcolor[rgb]{0.9,0.9,0.9}{
            \makecell[{{p{\linewidth}}}]{
                \texttt{\tiny{[GM$|$GM]}}
                \texttt{HAMLE: S, (3, 5)} \\
            }
        }
    }
    & & \\ \\

    \theutterance \stepcounter{utterance}  
    & & & \multicolumn{2}{p{0.3\linewidth}}{
        \cellcolor[rgb]{0.9,0.9,0.9}{
            \makecell[{{p{\linewidth}}}]{
                \texttt{\tiny{[GM$|$GM]}}
                \texttt{S moved to (3, 5)} \\
            }
        }
    }
    & & \\ \\

    \theutterance \stepcounter{utterance}  
    & & \multicolumn{4}{p{0.6\linewidth}}{
        \cellcolor[rgb]{0.9,0.9,0.9}{
            \makecell[{{p{\linewidth}}}]{
                \texttt{\tiny{[GM$\rangle$P2]}}
                \texttt{Mesajın diğer oyuncuya iletildi.} \\
\\ 
\\ 
\texttt{Şu anda en fazla 28 turdan 16. turu oynuyorsunuz.} \\
\texttt{Toplamda 16 cezanın 4 kadarını biriktirdiniz.Diğer oyuncu kendi ızgarasında bir nesneyi hareket ettirdi.} \\
\\ 
\\ 
\texttt{Bir sonraki komutun nedir?} \\
            }
        }
    }
    & \\ \\

    \theutterance \stepcounter{utterance}  
    & & & \multicolumn{4}{p{0.6\linewidth}}{
        \cellcolor[rgb]{0.9,0.9,0.9}{
            \makecell[{{p{\linewidth}}}]{
                \texttt{\tiny{[GM$\langle$P2]}}
                \texttt{HAMLE: P, (2,5)} \\
            }
        }
    }
     \\ \\

    \theutterance \stepcounter{utterance}  
    & & & \multicolumn{2}{p{0.3\linewidth}}{
        \cellcolor[rgb]{0.9,0.9,0.9}{
            \makecell[{{p{\linewidth}}}]{
                \texttt{\tiny{[GM$|$GM]}}
                \texttt{HAMLE: P, (2,5)} \\
            }
        }
    }
    & & \\ \\

    \theutterance \stepcounter{utterance}  
    & & & \multicolumn{2}{p{0.3\linewidth}}{
        \cellcolor[rgb]{0.9,0.9,0.9}{
            \makecell[{{p{\linewidth}}}]{
                \texttt{\tiny{[GM$|$GM]}}
                \texttt{P moved to (2, 5)} \\
            }
        }
    }
    & & \\ \\

    \theutterance \stepcounter{utterance}  
    & & \multicolumn{4}{p{0.6\linewidth}}{
        \cellcolor[rgb]{0.9,0.9,0.9}{
            \makecell[{{p{\linewidth}}}]{
                \texttt{\tiny{[P1$\langle$GM]}}
                \texttt{S’ başarıyla (3,5) koordinatlarına taşındı. Güncellenmiş ızgaran şu şekilde:} \\
\texttt{```} \\
\texttt{ 1234567} \\
\texttt{╔══╤══╤═╗} \\
\texttt{║◌◌│◌M│◌║ 1} \\
\texttt{╟──┼──┴─╢ 2} \\
\texttt{║◌◌│◌◌◌O║ 3} \\
\texttt{╟──┘◌A◌◌║ 4} \\
\texttt{║DPS◌◌◌◌║ 5} \\
\texttt{║◌◌┌─┬──╢ 6} \\
\texttt{║T◌│◌│◌◌║ 7} \\
\texttt{╚══╧═╧══╝} \\
\texttt{```} \\
\texttt{Şu anda en fazla 28 turdan 16. turu oynuyorsunuz.} \\
\texttt{Toplamda 16 cezanın 4 kadarını biriktirdiniz.Diğer oyuncu kendi ızgarasında bir nesneyi hareket ettirdi.} \\
\\ 
\\ 
\texttt{Bir sonraki komutun nedir?} \\
            }
        }
    }
    & \\ \\

    \theutterance \stepcounter{utterance}  
    & \multicolumn{4}{p{0.6\linewidth}}{
        \cellcolor[rgb]{0.9,0.9,0.9}{
            \makecell[{{p{\linewidth}}}]{
                \texttt{\tiny{[P1$\rangle$GM]}}
                \texttt{SÖYLE: S'yi (3,5)'e taşıdım. P(2,5) ve S(3,5) tamam. Şimdi M'yi (2,1)'den (4,5)'e taşıyacağım.} \\
            }
        }
    }
    & & \\ \\

    \theutterance \stepcounter{utterance}  
    & & & \multicolumn{2}{p{0.3\linewidth}}{
        \cellcolor[rgb]{0.9,0.9,0.9}{
            \makecell[{{p{\linewidth}}}]{
                \texttt{\tiny{[GM$|$GM]}}
                \texttt{SÖYLE: S'yi (3,5)'e taşıdım. P(2,5) ve S(3,5) tamam. Şimdi M'yi (2,1)'den (4,5)'e taşıyacağım.} \\
            }
        }
    }
    & & \\ \\

    \theutterance \stepcounter{utterance}  
    & & \multicolumn{4}{p{0.6\linewidth}}{
        \cellcolor[rgb]{0.9,0.9,0.9}{
            \makecell[{{p{\linewidth}}}]{
                \texttt{\tiny{[GM$\rangle$P2]}}
                \texttt{P’ başarıyla (2,5) koordinatlarına taşındı. Güncellenmiş ızgaran şu şekilde:} \\
\texttt{```} \\
\texttt{ 1234567} \\
\texttt{╔══╤══╤═╗} \\
\texttt{║◌◌│◌◌│T║ 1} \\
\texttt{╟──┼──┴─╢ 2} \\
\texttt{║◌◌│◌◌◌◌║ 3} \\
\texttt{╟──┘◌◌◌O║ 4} \\
\texttt{║DPSM◌◌◌║ 5} \\
\texttt{║◌◌┌─┬──╢ 6} \\
\texttt{║◌◌│A│◌◌║ 7} \\
\texttt{╚══╧═╧══╝} \\
\texttt{```} \\
\texttt{Şu anda en fazla 28 turdan 17. turu oynuyorsunuz.} \\
\texttt{Toplamda 16 cezanın 4 kadarını biriktirdiniz.Diğer oyuncu şu mesajı gönderdi:} \\
\\ 
\\ 
\texttt{"S'yi (3,5)'e taşıdım. P(2,5) ve S(3,5) tamam. Şimdi M'yi (2,1)'den (4,5)'e taşıyacağım."} \\
\\ 
\\ 
\texttt{Bir sonraki komutun nedir?} \\
            }
        }
    }
    & \\ \\

    \theutterance \stepcounter{utterance}  
    & & & \multicolumn{4}{p{0.6\linewidth}}{
        \cellcolor[rgb]{0.9,0.9,0.9}{
            \makecell[{{p{\linewidth}}}]{
                \texttt{\tiny{[GM$\langle$P2]}}
                \texttt{SÖYLE: Ben P(2,5), S(3,5), M(4,5) tamamladım! Şimdi O'yu (4,4)'ten (7,5)'e taşıyorum.} \\
            }
        }
    }
     \\ \\

    \theutterance \stepcounter{utterance}  
    & & & \multicolumn{2}{p{0.3\linewidth}}{
        \cellcolor[rgb]{0.9,0.9,0.9}{
            \makecell[{{p{\linewidth}}}]{
                \texttt{\tiny{[GM$|$GM]}}
                \texttt{SÖYLE: Ben P(2,5), S(3,5), M(4,5) tamamladım! Şimdi O'yu (4,4)'ten (7,5)'e taşıyorum.} \\
            }
        }
    }
    & & \\ \\

    \theutterance \stepcounter{utterance}  
    & & \multicolumn{4}{p{0.6\linewidth}}{
        \cellcolor[rgb]{0.9,0.9,0.9}{
            \makecell[{{p{\linewidth}}}]{
                \texttt{\tiny{[P1$\langle$GM]}}
                \texttt{Mesajın diğer oyuncuya iletildi.} \\
\\ 
\\ 
\texttt{Şu anda en fazla 28 turdan 17. turu oynuyorsunuz.} \\
\texttt{Toplamda 16 cezanın 4 kadarını biriktirdiniz.Diğer oyuncu şu mesajı gönderdi:} \\
\\ 
\\ 
\texttt{"Ben P(2,5), S(3,5), M(4,5) tamamladım! Şimdi O'yu (4,4)'ten (7,5)'e taşıyorum."} \\
\\ 
\\ 
\texttt{Bir sonraki komutun nedir?} \\
            }
        }
    }
    & \\ \\

    \theutterance \stepcounter{utterance}  
    & \multicolumn{4}{p{0.6\linewidth}}{
        \cellcolor[rgb]{0.9,0.9,0.9}{
            \makecell[{{p{\linewidth}}}]{
                \texttt{\tiny{[P1$\rangle$GM]}}
                \texttt{HAMLE: M, (4, 5)} \\
            }
        }
    }
    & & \\ \\

    \theutterance \stepcounter{utterance}  
    & & & \multicolumn{2}{p{0.3\linewidth}}{
        \cellcolor[rgb]{0.9,0.9,0.9}{
            \makecell[{{p{\linewidth}}}]{
                \texttt{\tiny{[GM$|$GM]}}
                \texttt{HAMLE: M, (4, 5)} \\
            }
        }
    }
    & & \\ \\

    \theutterance \stepcounter{utterance}  
    & & & \multicolumn{2}{p{0.3\linewidth}}{
        \cellcolor[rgb]{0.9,0.9,0.9}{
            \makecell[{{p{\linewidth}}}]{
                \texttt{\tiny{[GM$|$GM]}}
                \texttt{M moved to (4, 5)} \\
            }
        }
    }
    & & \\ \\

    \theutterance \stepcounter{utterance}  
    & & \multicolumn{4}{p{0.6\linewidth}}{
        \cellcolor[rgb]{0.9,0.9,0.9}{
            \makecell[{{p{\linewidth}}}]{
                \texttt{\tiny{[GM$\rangle$P2]}}
                \texttt{Mesajın diğer oyuncuya iletildi.} \\
\\ 
\\ 
\texttt{Şu anda en fazla 28 turdan 18. turu oynuyorsunuz.} \\
\texttt{Toplamda 16 cezanın 4 kadarını biriktirdiniz.Diğer oyuncu kendi ızgarasında bir nesneyi hareket ettirdi.} \\
\\ 
\\ 
\texttt{Bir sonraki komutun nedir?} \\
            }
        }
    }
    & \\ \\

    \theutterance \stepcounter{utterance}  
    & & & \multicolumn{4}{p{0.6\linewidth}}{
        \cellcolor[rgb]{0.9,0.9,0.9}{
            \makecell[{{p{\linewidth}}}]{
                \texttt{\tiny{[GM$\langle$P2]}}
                \texttt{HAMLE: O, (7,5)} \\
            }
        }
    }
     \\ \\

    \theutterance \stepcounter{utterance}  
    & & & \multicolumn{2}{p{0.3\linewidth}}{
        \cellcolor[rgb]{0.9,0.9,0.9}{
            \makecell[{{p{\linewidth}}}]{
                \texttt{\tiny{[GM$|$GM]}}
                \texttt{HAMLE: O, (7,5)} \\
            }
        }
    }
    & & \\ \\

    \theutterance \stepcounter{utterance}  
    & & & \multicolumn{2}{p{0.3\linewidth}}{
        \cellcolor[rgb]{0.9,0.9,0.9}{
            \makecell[{{p{\linewidth}}}]{
                \texttt{\tiny{[GM$|$GM]}}
                \texttt{O moved to (7, 5)} \\
            }
        }
    }
    & & \\ \\

    \theutterance \stepcounter{utterance}  
    & & \multicolumn{4}{p{0.6\linewidth}}{
        \cellcolor[rgb]{0.9,0.9,0.9}{
            \makecell[{{p{\linewidth}}}]{
                \texttt{\tiny{[P1$\langle$GM]}}
                \texttt{M’ başarıyla (4,5) koordinatlarına taşındı. Güncellenmiş ızgaran şu şekilde:} \\
\texttt{```} \\
\texttt{ 1234567} \\
\texttt{╔══╤══╤═╗} \\
\texttt{║◌◌│◌◌│◌║ 1} \\
\texttt{╟──┼──┴─╢ 2} \\
\texttt{║◌◌│◌◌◌O║ 3} \\
\texttt{╟──┘◌A◌◌║ 4} \\
\texttt{║DPSM◌◌◌║ 5} \\
\texttt{║◌◌┌─┬──╢ 6} \\
\texttt{║T◌│◌│◌◌║ 7} \\
\texttt{╚══╧═╧══╝} \\
\texttt{```} \\
\texttt{Şu anda en fazla 28 turdan 18. turu oynuyorsunuz.} \\
\texttt{Toplamda 16 cezanın 4 kadarını biriktirdiniz.Diğer oyuncu kendi ızgarasında bir nesneyi hareket ettirdi.} \\
\\ 
\\ 
\texttt{Bir sonraki komutun nedir?} \\
            }
        }
    }
    & \\ \\

    \theutterance \stepcounter{utterance}  
    & \multicolumn{4}{p{0.6\linewidth}}{
        \cellcolor[rgb]{0.9,0.9,0.9}{
            \makecell[{{p{\linewidth}}}]{
                \texttt{\tiny{[P1$\rangle$GM]}}
                \texttt{SÖYLE: M'yi (4,5)'e taşıdım! P(2,5), S(3,5), M(4,5) tamam. Şimdi O'yu (7,3)'ten (7,5)'e taşıyacağım. D(1,5) zaten yerinde. T'yi (1,7)'den (7,1)'e taşımam lazım. A'nın durumu ne? Sen A'yı nereye taşıdın?} \\
            }
        }
    }
    & & \\ \\

    \theutterance \stepcounter{utterance}  
    & & & \multicolumn{2}{p{0.3\linewidth}}{
        \cellcolor[rgb]{0.9,0.9,0.9}{
            \makecell[{{p{\linewidth}}}]{
                \texttt{\tiny{[GM$|$GM]}}
                \texttt{SÖYLE: M'yi (4,5)'e taşıdım! P(2,5), S(3,5), M(4,5) tamam. Şimdi O'yu (7,3)'ten (7,5)'e taşıyacağım. D(1,5) zaten yerinde. T'yi (1,7)'den (7,1)'e taşımam lazım. A'nın durumu ne? Sen A'yı nereye taşıdın?} \\
            }
        }
    }
    & & \\ \\

    \theutterance \stepcounter{utterance}  
    & & \multicolumn{4}{p{0.6\linewidth}}{
        \cellcolor[rgb]{0.9,0.9,0.9}{
            \makecell[{{p{\linewidth}}}]{
                \texttt{\tiny{[GM$\rangle$P2]}}
                \texttt{O’ başarıyla (7,5) koordinatlarına taşındı. Güncellenmiş ızgaran şu şekilde:} \\
\texttt{```} \\
\texttt{ 1234567} \\
\texttt{╔══╤══╤═╗} \\
\texttt{║◌◌│◌◌│T║ 1} \\
\texttt{╟──┼──┴─╢ 2} \\
\texttt{║◌◌│◌◌◌◌║ 3} \\
\texttt{╟──┘◌◌◌◌║ 4} \\
\texttt{║DPSM◌◌O║ 5} \\
\texttt{║◌◌┌─┬──╢ 6} \\
\texttt{║◌◌│A│◌◌║ 7} \\
\texttt{╚══╧═╧══╝} \\
\texttt{```} \\
\texttt{Şu anda en fazla 28 turdan 19. turu oynuyorsunuz.} \\
\texttt{Toplamda 16 cezanın 4 kadarını biriktirdiniz.Diğer oyuncu şu mesajı gönderdi:} \\
\\ 
\\ 
\texttt{"M'yi (4,5)'e taşıdım! P(2,5), S(3,5), M(4,5) tamam. Şimdi O'yu (7,3)'ten (7,5)'e taşıyacağım. D(1,5) zaten yerinde. T'yi (1,7)'den (7,1)'e taşımam lazım. A'nın durumu ne? Sen A'yı nereye taşıdın?"} \\
\\ 
\\ 
\texttt{Bir sonraki komutun nedir?} \\
            }
        }
    }
    & \\ \\

    \theutterance \stepcounter{utterance}  
    & & & \multicolumn{4}{p{0.6\linewidth}}{
        \cellcolor[rgb]{0.9,0.9,0.9}{
            \makecell[{{p{\linewidth}}}]{
                \texttt{\tiny{[GM$\langle$P2]}}
                \texttt{SÖYLE: O'yu (7,5)'e taşıdım, tamam! A(3,7) olduğu yerde kalıyor—tahtamda (3,7) dışında A için boş yer bulmak zor, çizgilerle dolu. T'yi (3,1)'den (7,1)'e taşımaya çalışıyorum şimdi.} \\
            }
        }
    }
     \\ \\

    \theutterance \stepcounter{utterance}  
    & & & \multicolumn{2}{p{0.3\linewidth}}{
        \cellcolor[rgb]{0.9,0.9,0.9}{
            \makecell[{{p{\linewidth}}}]{
                \texttt{\tiny{[GM$|$GM]}}
                \texttt{SÖYLE: O'yu (7,5)'e taşıdım, tamam! A(3,7) olduğu yerde kalıyor—tahtamda (3,7) dışında A için boş yer bulmak zor, çizgilerle dolu. T'yi (3,1)'den (7,1)'e taşımaya çalışıyorum şimdi.} \\
            }
        }
    }
    & & \\ \\

    \theutterance \stepcounter{utterance}  
    & & \multicolumn{4}{p{0.6\linewidth}}{
        \cellcolor[rgb]{0.9,0.9,0.9}{
            \makecell[{{p{\linewidth}}}]{
                \texttt{\tiny{[P1$\langle$GM]}}
                \texttt{Mesajın diğer oyuncuya iletildi.} \\
\\ 
\\ 
\texttt{Şu anda en fazla 28 turdan 19. turu oynuyorsunuz.} \\
\texttt{Toplamda 16 cezanın 4 kadarını biriktirdiniz.Diğer oyuncu şu mesajı gönderdi:} \\
\\ 
\\ 
\texttt{"O'yu (7,5)'e taşıdım, tamam! A(3,7) olduğu yerde kalıyor—tahtamda (3,7) dışında A için boş yer bulmak zor, çizgilerle dolu. T'yi (3,1)'den (7,1)'e taşımaya çalışıyorum şimdi."} \\
\\ 
\\ 
\texttt{Bir sonraki komutun nedir?} \\
            }
        }
    }
    & \\ \\

    \theutterance \stepcounter{utterance}  
    & \multicolumn{4}{p{0.6\linewidth}}{
        \cellcolor[rgb]{0.9,0.9,0.9}{
            \makecell[{{p{\linewidth}}}]{
                \texttt{\tiny{[P1$\rangle$GM]}}
                \texttt{HAMLE: O, (7, 5)} \\
            }
        }
    }
    & & \\ \\

    \theutterance \stepcounter{utterance}  
    & & & \multicolumn{2}{p{0.3\linewidth}}{
        \cellcolor[rgb]{0.9,0.9,0.9}{
            \makecell[{{p{\linewidth}}}]{
                \texttt{\tiny{[GM$|$GM]}}
                \texttt{HAMLE: O, (7, 5)} \\
            }
        }
    }
    & & \\ \\

    \theutterance \stepcounter{utterance}  
    & & & \multicolumn{2}{p{0.3\linewidth}}{
        \cellcolor[rgb]{0.9,0.9,0.9}{
            \makecell[{{p{\linewidth}}}]{
                \texttt{\tiny{[GM$|$GM]}}
                \texttt{O moved to (7, 5)} \\
            }
        }
    }
    & & \\ \\

    \theutterance \stepcounter{utterance}  
    & & \multicolumn{4}{p{0.6\linewidth}}{
        \cellcolor[rgb]{0.9,0.9,0.9}{
            \makecell[{{p{\linewidth}}}]{
                \texttt{\tiny{[GM$\rangle$P2]}}
                \texttt{Mesajın diğer oyuncuya iletildi.} \\
\\ 
\\ 
\texttt{Şu anda en fazla 28 turdan 20. turu oynuyorsunuz.} \\
\texttt{Toplamda 16 cezanın 4 kadarını biriktirdiniz.Diğer oyuncu kendi ızgarasında bir nesneyi hareket ettirdi.} \\
\\ 
\\ 
\texttt{Bir sonraki komutun nedir?} \\
            }
        }
    }
    & \\ \\

    \theutterance \stepcounter{utterance}  
    & & & \multicolumn{4}{p{0.6\linewidth}}{
        \cellcolor[rgb]{0.9,0.9,0.9}{
            \makecell[{{p{\linewidth}}}]{
                \texttt{\tiny{[GM$\langle$P2]}}
                \texttt{HAMLE: T, (7,1)} \\
            }
        }
    }
     \\ \\

    \theutterance \stepcounter{utterance}  
    & & & \multicolumn{2}{p{0.3\linewidth}}{
        \cellcolor[rgb]{0.9,0.9,0.9}{
            \makecell[{{p{\linewidth}}}]{
                \texttt{\tiny{[GM$|$GM]}}
                \texttt{HAMLE: T, (7,1)} \\
            }
        }
    }
    & & \\ \\

    \theutterance \stepcounter{utterance}  
    & & & \multicolumn{2}{p{0.3\linewidth}}{
        \cellcolor[rgb]{0.9,0.9,0.9}{
            \makecell[{{p{\linewidth}}}]{
                \texttt{\tiny{[GM$|$GM]}}
                \texttt{Ceza: (7,1) boş değil; "T" içeriyor.} \\
\\ 
\texttt{Toplamda 16 cezanın 5 kadarını biriktirdiniz.Lütfen tekrar deneyin!} \\
            }
        }
    }
    & & \\ \\

    \theutterance \stepcounter{utterance}  
    & & \multicolumn{4}{p{0.6\linewidth}}{
        \cellcolor[rgb]{0.9,0.9,0.9}{
            \makecell[{{p{\linewidth}}}]{
                \texttt{\tiny{[GM$\rangle$P2]}}
                \texttt{Ceza: (7,1) boş değil; "T" içeriyor.} \\
\\ 
\texttt{Toplamda 16 cezanın 5 kadarını biriktirdiniz.Lütfen tekrar deneyin!} \\
            }
        }
    }
    & \\ \\

    \theutterance \stepcounter{utterance}  
    & & & \multicolumn{4}{p{0.6\linewidth}}{
        \cellcolor[rgb]{0.9,0.9,0.9}{
            \makecell[{{p{\linewidth}}}]{
                \texttt{\tiny{[GM$\langle$P2]}}
                \texttt{SÖYLE: T zaten (7,1)'de! Yani T'ye dokunmama gerek yok. Durumum: P(2,5), S(3,5), M(4,5), O(7,5), T(7,1), D(6,5), A(3,7) — hepsi yerinde! Senin durum ne? T(7,1) ve O(7,5) tamam mı?} \\
            }
        }
    }
     \\ \\

    \theutterance \stepcounter{utterance}  
    & & & \multicolumn{2}{p{0.3\linewidth}}{
        \cellcolor[rgb]{0.9,0.9,0.9}{
            \makecell[{{p{\linewidth}}}]{
                \texttt{\tiny{[GM$|$GM]}}
                \texttt{SÖYLE: T zaten (7,1)'de! Yani T'ye dokunmama gerek yok. Durumum: P(2,5), S(3,5), M(4,5), O(7,5), T(7,1), D(6,5), A(3,7) — hepsi yerinde! Senin durum ne? T(7,1) ve O(7,5) tamam mı?} \\
            }
        }
    }
    & & \\ \\

    \theutterance \stepcounter{utterance}  
    & & \multicolumn{4}{p{0.6\linewidth}}{
        \cellcolor[rgb]{0.9,0.9,0.9}{
            \makecell[{{p{\linewidth}}}]{
                \texttt{\tiny{[P1$\langle$GM]}}
                \texttt{O’ başarıyla (7,5) koordinatlarına taşındı. Güncellenmiş ızgaran şu şekilde:} \\
\texttt{```} \\
\texttt{ 1234567} \\
\texttt{╔══╤══╤═╗} \\
\texttt{║◌◌│◌◌│◌║ 1} \\
\texttt{╟──┼──┴─╢ 2} \\
\texttt{║◌◌│◌◌◌◌║ 3} \\
\texttt{╟──┘◌A◌◌║ 4} \\
\texttt{║DPSM◌◌O║ 5} \\
\texttt{║◌◌┌─┬──╢ 6} \\
\texttt{║T◌│◌│◌◌║ 7} \\
\texttt{╚══╧═╧══╝} \\
\texttt{```} \\
\texttt{Şu anda en fazla 28 turdan 20. turu oynuyorsunuz.} \\
\texttt{Toplamda 16 cezanın 5 kadarını biriktirdiniz.Diğer oyuncu şu mesajı gönderdi:} \\
\\ 
\\ 
\texttt{"T zaten (7,1)'de! Yani T'ye dokunmama gerek yok. Durumum: P(2,5), S(3,5), M(4,5), O(7,5), T(7,1), D(6,5), A(3,7) — hepsi yerinde! Senin durum ne? T(7,1) ve O(7,5) tamam mı?"} \\
\\ 
\\ 
\texttt{Bir sonraki komutun nedir?} \\
            }
        }
    }
    & \\ \\

    \theutterance \stepcounter{utterance}  
    & \multicolumn{4}{p{0.6\linewidth}}{
        \cellcolor[rgb]{0.9,0.9,0.9}{
            \makecell[{{p{\linewidth}}}]{
                \texttt{\tiny{[P1$\rangle$GM]}}
                \texttt{SÖYLE: O(7,5) tamam! Şimdi T'yi (1,7)'den (7,1)'e taşıyorum. D(1,5) zaten yerinde. A(3,4)'e yakın bir yere taşıyabilir miyim? Senin A(3,7)'de, benim A(3,4)'te — yakın ama aynı değil. A için (3,7) senin tahtanda uygun ama benim (3,7) dolu. A'yı olduğu yerde bırakalım, skoru etkileyecek ama T'yi önce taşıyayım.} \\
            }
        }
    }
    & & \\ \\

    \theutterance \stepcounter{utterance}  
    & & & \multicolumn{2}{p{0.3\linewidth}}{
        \cellcolor[rgb]{0.9,0.9,0.9}{
            \makecell[{{p{\linewidth}}}]{
                \texttt{\tiny{[GM$|$GM]}}
                \texttt{SÖYLE: O(7,5) tamam! Şimdi T'yi (1,7)'den (7,1)'e taşıyorum. D(1,5) zaten yerinde. A(3,4)'e yakın bir yere taşıyabilir miyim? Senin A(3,7)'de, benim A(3,4)'te — yakın ama aynı değil. A için (3,7) senin tahtanda uygun ama benim (3,7) dolu. A'yı olduğu yerde bırakalım, skoru etkileyecek ama T'yi önce taşıyayım.} \\
            }
        }
    }
    & & \\ \\

    \theutterance \stepcounter{utterance}  
    & & \multicolumn{4}{p{0.6\linewidth}}{
        \cellcolor[rgb]{0.9,0.9,0.9}{
            \makecell[{{p{\linewidth}}}]{
                \texttt{\tiny{[GM$\rangle$P2]}}
                \texttt{Mesajın diğer oyuncuya iletildi.} \\
\\ 
\\ 
\texttt{Şu anda en fazla 28 turdan 21. turu oynuyorsunuz.} \\
\texttt{Toplamda 16 cezanın 5 kadarını biriktirdiniz.Diğer oyuncu şu mesajı gönderdi:} \\
\\ 
\\ 
\texttt{"O(7,5) tamam! Şimdi T'yi (1,7)'den (7,1)'e taşıyorum. D(1,5) zaten yerinde. A(3,4)'e yakın bir yere taşıyabilir miyim? Senin A(3,7)'de, benim A(3,4)'te — yakın ama aynı değil. A için (3,7) senin tahtanda uygun ama benim (3,7) dolu. A'yı olduğu yerde bırakalım, skoru etkileyecek ama T'yi önce taşıyayım."} \\
\\ 
\\ 
\texttt{Bir sonraki komutun nedir?} \\
            }
        }
    }
    & \\ \\

    \theutterance \stepcounter{utterance}  
    & & & \multicolumn{4}{p{0.6\linewidth}}{
        \cellcolor[rgb]{0.9,0.9,0.9}{
            \makecell[{{p{\linewidth}}}]{
                \texttt{\tiny{[GM$\langle$P2]}}
                \texttt{SÖYLE: Benim tüm nesnelerim yerinde: P(2,5), S(3,5), M(4,5), O(7,5), T(7,1), D(6,5), A(3,7). Sen T'yi taşıdıktan sonra bitti mi? diye sorabilirsin. A'yı olduğu yerde bırakıyoruz—başka çare yok gibi.} \\
            }
        }
    }
     \\ \\

    \theutterance \stepcounter{utterance}  
    & & & \multicolumn{2}{p{0.3\linewidth}}{
        \cellcolor[rgb]{0.9,0.9,0.9}{
            \makecell[{{p{\linewidth}}}]{
                \texttt{\tiny{[GM$|$GM]}}
                \texttt{SÖYLE: Benim tüm nesnelerim yerinde: P(2,5), S(3,5), M(4,5), O(7,5), T(7,1), D(6,5), A(3,7). Sen T'yi taşıdıktan sonra bitti mi? diye sorabilirsin. A'yı olduğu yerde bırakıyoruz—başka çare yok gibi.} \\
            }
        }
    }
    & & \\ \\

    \theutterance \stepcounter{utterance}  
    & & \multicolumn{4}{p{0.6\linewidth}}{
        \cellcolor[rgb]{0.9,0.9,0.9}{
            \makecell[{{p{\linewidth}}}]{
                \texttt{\tiny{[P1$\langle$GM]}}
                \texttt{Mesajın diğer oyuncuya iletildi.} \\
\\ 
\\ 
\texttt{Şu anda en fazla 28 turdan 21. turu oynuyorsunuz.} \\
\texttt{Toplamda 16 cezanın 5 kadarını biriktirdiniz.Diğer oyuncu şu mesajı gönderdi:} \\
\\ 
\\ 
\texttt{"Benim tüm nesnelerim yerinde: P(2,5), S(3,5), M(4,5), O(7,5), T(7,1), D(6,5), A(3,7). Sen T'yi taşıdıktan sonra bitti mi? diye sorabilirsin. A'yı olduğu yerde bırakıyoruz—başka çare yok gibi."} \\
\\ 
\\ 
\texttt{Bir sonraki komutun nedir?} \\
            }
        }
    }
    & \\ \\

    \theutterance \stepcounter{utterance}  
    & \multicolumn{4}{p{0.6\linewidth}}{
        \cellcolor[rgb]{0.9,0.9,0.9}{
            \makecell[{{p{\linewidth}}}]{
                \texttt{\tiny{[P1$\rangle$GM]}}
                \texttt{HAMLE: T, (7, 1)} \\
            }
        }
    }
    & & \\ \\

    \theutterance \stepcounter{utterance}  
    & & & \multicolumn{2}{p{0.3\linewidth}}{
        \cellcolor[rgb]{0.9,0.9,0.9}{
            \makecell[{{p{\linewidth}}}]{
                \texttt{\tiny{[GM$|$GM]}}
                \texttt{HAMLE: T, (7, 1)} \\
            }
        }
    }
    & & \\ \\

    \theutterance \stepcounter{utterance}  
    & & & \multicolumn{2}{p{0.3\linewidth}}{
        \cellcolor[rgb]{0.9,0.9,0.9}{
            \makecell[{{p{\linewidth}}}]{
                \texttt{\tiny{[GM$|$GM]}}
                \texttt{T moved to (7, 1)} \\
            }
        }
    }
    & & \\ \\

    \theutterance \stepcounter{utterance}  
    & & \multicolumn{4}{p{0.6\linewidth}}{
        \cellcolor[rgb]{0.9,0.9,0.9}{
            \makecell[{{p{\linewidth}}}]{
                \texttt{\tiny{[GM$\rangle$P2]}}
                \texttt{Mesajın diğer oyuncuya iletildi.} \\
\\ 
\\ 
\texttt{Şu anda en fazla 28 turdan 22. turu oynuyorsunuz.} \\
\texttt{Toplamda 16 cezanın 5 kadarını biriktirdiniz.Diğer oyuncu kendi ızgarasında bir nesneyi hareket ettirdi.} \\
\\ 
\\ 
\texttt{Bir sonraki komutun nedir?} \\
            }
        }
    }
    & \\ \\

    \theutterance \stepcounter{utterance}  
    & & & \multicolumn{4}{p{0.6\linewidth}}{
        \cellcolor[rgb]{0.9,0.9,0.9}{
            \makecell[{{p{\linewidth}}}]{
                \texttt{\tiny{[GM$\langle$P2]}}
                \texttt{SÖYLE: Benim tüm nesnelerim yerinde. T(7,1) taşıman tamam mı? Hazırsan "bitti mi?" diye soruyorum.} \\
            }
        }
    }
     \\ \\

    \theutterance \stepcounter{utterance}  
    & & & \multicolumn{2}{p{0.3\linewidth}}{
        \cellcolor[rgb]{0.9,0.9,0.9}{
            \makecell[{{p{\linewidth}}}]{
                \texttt{\tiny{[GM$|$GM]}}
                \texttt{SÖYLE: Benim tüm nesnelerim yerinde. T(7,1) taşıman tamam mı? Hazırsan "bitti mi?" diye soruyorum.} \\
            }
        }
    }
    & & \\ \\

    \theutterance \stepcounter{utterance}  
    & & \multicolumn{4}{p{0.6\linewidth}}{
        \cellcolor[rgb]{0.9,0.9,0.9}{
            \makecell[{{p{\linewidth}}}]{
                \texttt{\tiny{[P1$\langle$GM]}}
                \texttt{T’ başarıyla (7,1) koordinatlarına taşındı. Güncellenmiş ızgaran şu şekilde:} \\
\texttt{```} \\
\texttt{ 1234567} \\
\texttt{╔══╤══╤═╗} \\
\texttt{║◌◌│◌◌│T║ 1} \\
\texttt{╟──┼──┴─╢ 2} \\
\texttt{║◌◌│◌◌◌◌║ 3} \\
\texttt{╟──┘◌A◌◌║ 4} \\
\texttt{║DPSM◌◌O║ 5} \\
\texttt{║◌◌┌─┬──╢ 6} \\
\texttt{║◌◌│◌│◌◌║ 7} \\
\texttt{╚══╧═╧══╝} \\
\texttt{```} \\
\texttt{Şu anda en fazla 28 turdan 22. turu oynuyorsunuz.} \\
\texttt{Toplamda 16 cezanın 5 kadarını biriktirdiniz.Diğer oyuncu şu mesajı gönderdi:} \\
\\ 
\\ 
\texttt{"Benim tüm nesnelerim yerinde. T(7,1) taşıman tamam mı? Hazırsan "bitti mi?" diye soruyorum."} \\
\\ 
\\ 
\texttt{Bir sonraki komutun nedir?} \\
            }
        }
    }
    & \\ \\

    \theutterance \stepcounter{utterance}  
    & \multicolumn{4}{p{0.6\linewidth}}{
        \cellcolor[rgb]{0.9,0.9,0.9}{
            \makecell[{{p{\linewidth}}}]{
                \texttt{\tiny{[P1$\rangle$GM]}}
                \texttt{SÖYLE: T(7,1) tamam! Benim tüm nesnelerim yerinde: P(2,5), S(3,5), M(4,5), O(7,5), T(7,1), D(1,5), A(3,4). D hariç hepsi aynı konumda—D bende (1,5), sende (6,5). A bende (3,4), sende (3,7). Bitti!} \\
            }
        }
    }
    & & \\ \\

    \theutterance \stepcounter{utterance}  
    & & & \multicolumn{2}{p{0.3\linewidth}}{
        \cellcolor[rgb]{0.9,0.9,0.9}{
            \makecell[{{p{\linewidth}}}]{
                \texttt{\tiny{[GM$|$GM]}}
                \texttt{SÖYLE: T(7,1) tamam! Benim tüm nesnelerim yerinde: P(2,5), S(3,5), M(4,5), O(7,5), T(7,1), D(1,5), A(3,4). D hariç hepsi aynı konumda—D bende (1,5), sende (6,5). A bende (3,4), sende (3,7). Bitti!} \\
            }
        }
    }
    & & \\ \\

    \theutterance \stepcounter{utterance}  
    & & \multicolumn{4}{p{0.6\linewidth}}{
        \cellcolor[rgb]{0.9,0.9,0.9}{
            \makecell[{{p{\linewidth}}}]{
                \texttt{\tiny{[GM$\rangle$P2]}}
                \texttt{Mesajın diğer oyuncuya iletildi.} \\
\\ 
\\ 
\texttt{Şu anda en fazla 28 turdan 23. turu oynuyorsunuz.} \\
\texttt{Toplamda 16 cezanın 5 kadarını biriktirdiniz.Diğer oyuncu şu mesajı gönderdi:} \\
\\ 
\\ 
\texttt{"T(7,1) tamam! Benim tüm nesnelerim yerinde: P(2,5), S(3,5), M(4,5), O(7,5), T(7,1), D(1,5), A(3,4). D hariç hepsi aynı konumda—D bende (1,5), sende (6,5). A bende (3,4), sende (3,7). Bitti!"} \\
\\ 
\\ 
\texttt{Bir sonraki komutun nedir?} \\
            }
        }
    }
    & \\ \\

    \theutterance \stepcounter{utterance}  
    & & & \multicolumn{4}{p{0.6\linewidth}}{
        \cellcolor[rgb]{0.9,0.9,0.9}{
            \makecell[{{p{\linewidth}}}]{
                \texttt{\tiny{[GM$\langle$P2]}}
                \texttt{SÖYLE: bitti!} \\
            }
        }
    }
     \\ \\

    \theutterance \stepcounter{utterance}  
    & & & \multicolumn{2}{p{0.3\linewidth}}{
        \cellcolor[rgb]{0.9,0.9,0.9}{
            \makecell[{{p{\linewidth}}}]{
                \texttt{\tiny{[GM$|$GM]}}
                \texttt{SÖYLE: bitti!} \\
            }
        }
    }
    & & \\ \\

    \theutterance \stepcounter{utterance}  
    & & \multicolumn{4}{p{0.6\linewidth}}{
        \cellcolor[rgb]{0.9,0.9,0.9}{
            \makecell[{{p{\linewidth}}}]{
                \texttt{\tiny{[P1$\langle$GM]}}
                \texttt{Mesajın diğer oyuncuya iletildi.} \\
\\ 
\\ 
\texttt{Şu anda en fazla 28 turdan 23. turu oynuyorsunuz.} \\
\texttt{Toplamda 16 cezanın 5 kadarını biriktirdiniz.Diğer oyuncu şu mesajı gönderdi:} \\
\\ 
\\ 
\texttt{"bitti!"} \\
\\ 
\\ 
\texttt{Bir sonraki komutun nedir?} \\
            }
        }
    }
    & \\ \\

    \theutterance \stepcounter{utterance}  
    & \multicolumn{4}{p{0.6\linewidth}}{
        \cellcolor[rgb]{0.9,0.9,0.9}{
            \makecell[{{p{\linewidth}}}]{
                \texttt{\tiny{[P1$\rangle$GM]}}
                \texttt{SÖYLE: bitti!} \\
            }
        }
    }
    & & \\ \\

    \theutterance \stepcounter{utterance}  
    & & & \multicolumn{2}{p{0.3\linewidth}}{
        \cellcolor[rgb]{0.9,0.9,0.9}{
            \makecell[{{p{\linewidth}}}]{
                \texttt{\tiny{[GM$|$GM]}}
                \texttt{SÖYLE: bitti!} \\
            }
        }
    }
    & & \\ \\

    \theutterance \stepcounter{utterance}  
    & & & \multicolumn{2}{p{0.3\linewidth}}{
        \cellcolor[rgb]{0.9,0.9,0.9}{
            \makecell[{{p{\linewidth}}}]{
                \texttt{\tiny{[GM$|$GM]}}
                \texttt{true} \\
            }
        }
    }
    & & \\ \\

    \theutterance \stepcounter{utterance}  
    & & & \multicolumn{2}{p{0.3\linewidth}}{
        \cellcolor[rgb]{0.9,0.9,0.9}{
            \makecell[{{p{\linewidth}}}]{
                \texttt{\tiny{[GM$|$GM]}}
                \texttt{Initial states:} \\
\texttt{```} \\
\texttt{Player 1:} \\
\texttt{ 1234567} \\
\texttt{╔══╤══╤═╗} \\
\texttt{║◌P│◌M│◌║ 1} \\
\texttt{╟──┼──┴─╢ 2} \\
\texttt{║◌◌│◌S◌O║ 3} \\
\texttt{╟──┘◌A◌◌║ 4} \\
\texttt{║D◌◌◌◌◌◌║ 5} \\
\texttt{║◌◌┌─┬──╢ 6} \\
\texttt{║T◌│◌│◌◌║ 7} \\
\texttt{╚══╧═╧══╝} \\
\\ 
\\ 
\texttt{Player 2:} \\
\texttt{ 1234567} \\
\texttt{╔══╤══╤═╗} \\
\texttt{║◌◌│◌◌│T║ 1} \\
\texttt{╟──┼──┴─╢ 2} \\
\texttt{║◌◌│◌S◌P║ 3} \\
\texttt{╟──┘◌◌MO║ 4} \\
\texttt{║D◌◌◌◌◌◌║ 5} \\
\texttt{║◌◌┌─┬──╢ 6} \\
\texttt{║◌◌│A│◌◌║ 7} \\
\texttt{╚══╧═╧══╝} \\
\\ 
\texttt{```} \\
            }
        }
    }
    & & \\ \\

    \theutterance \stepcounter{utterance}  
    & & & \multicolumn{2}{p{0.3\linewidth}}{
        \cellcolor[rgb]{0.9,0.9,0.9}{
            \makecell[{{p{\linewidth}}}]{
                \texttt{\tiny{[GM$|$GM]}}
                \texttt{End states:} \\
\texttt{```} \\
\texttt{Player 1:} \\
\texttt{ 1234567} \\
\texttt{╔══╤══╤═╗} \\
\texttt{║◌◌│◌◌│T║ 1} \\
\texttt{╟──┼──┴─╢ 2} \\
\texttt{║◌◌│◌◌◌◌║ 3} \\
\texttt{╟──┘◌A◌◌║ 4} \\
\texttt{║DPSM◌◌O║ 5} \\
\texttt{║◌◌┌─┬──╢ 6} \\
\texttt{║◌◌│◌│◌◌║ 7} \\
\texttt{╚══╧═╧══╝} \\
\\ 
\\ 
\texttt{Player 2:} \\
\texttt{ 1234567} \\
\texttt{╔══╤══╤═╗} \\
\texttt{║◌◌│◌◌│T║ 1} \\
\texttt{╟──┼──┴─╢ 2} \\
\texttt{║◌◌│◌◌◌◌║ 3} \\
\texttt{╟──┘◌◌◌◌║ 4} \\
\texttt{║DPSM◌◌O║ 5} \\
\texttt{║◌◌┌─┬──╢ 6} \\
\texttt{║◌◌│A│◌◌║ 7} \\
\texttt{╚══╧═╧══╝} \\
\\ 
\texttt{```} \\
            }
        }
    }
    & & \\ \\

    \theutterance \stepcounter{utterance}  
    & & & \multicolumn{2}{p{0.3\linewidth}}{
        \cellcolor[rgb]{0.9,0.9,0.9}{
            \makecell[{{p{\linewidth}}}]{
                \texttt{\tiny{[GM$|$GM]}}
                \texttt{* success: True} \\
\texttt{* lose: False} \\
\texttt{* aborted: False} \\
\texttt{{-}{-}{-}{-}{-}{-}{-}} \\
\texttt{* Shifts: 7.00} \\
\texttt{* Max Shifts: 12.00} \\
\texttt{* Min Shifts: 6.00} \\
\texttt{* End Distance Sum: 3.16} \\
\texttt{* Init Distance Sum: 21.20} \\
\texttt{* Expected Distance Sum: 29.33} \\
\texttt{* Penalties: 5.00} \\
\texttt{* Max Penalties: 16.00} \\
\texttt{* Rounds: 23.00} \\
\texttt{* Max Rounds: 28.00} \\
\texttt{* Object Count: 7.00} \\
            }
        }
    }
    & & \\ \\

\end{supertabular}
}

\end{document}
